% ---------------------------------------------------------------
% Modelo LaTex para dissertação e tese do programa de Pós Graduação em Ciência da Computação da UFABC
% ---------------------------------------------------------------

\documentclass[
	% -- opções da classe memoir --
	12pt,					% tamanho da fonte
	openright,				% capítulos começam em pág ímpar (insere página vazia caso preciso)
	twoside,					% para impressão em verso e anverso. Oposto a oneside
	a4paper,					% tamanho do papel. 
	% -- opções da classe abntex2 --
	%chapter=TITLE,			% títulos de capítulos convertidos em letras maiúsculas
	%section=TITLE,			% títulos de seções convertidos em letras maiúsculas
	%subsection=TITLE,		% títulos de subseções convertidos em letras maiúsculas
	%subsubsection=TITLE,	% títulos de subsubseções convertidos em letras maiúsculas
	% -- opções do pacote babel --
	english,					% idioma adicional para hifenização
	%french,					% idioma adicional para hifenização
	%spanish,				% idioma adicional para hifenização
	brazil					% o último idioma é o principal do documento
	]{abntex2}

% ---------------------
% Pacotes OBRIGATÓRIOS
% ---------------------
\usepackage{lmodern}				% Usa a fonte Latin Modern			
\usepackage[T1]{fontenc}			% Selecao de codigos de fonte.
\usepackage[utf8]{inputenc}		% Codificacao do documento (conversão automática dos acentos)
\usepackage{textcomp}
\usepackage{lastpage}			% Usado pela Ficha catalográfica
\usepackage{indentfirst}			% Indenta o primeiro parágrafo de cada seção.
\usepackage{color}				% Controle das cores
\usepackage{graphicx,graphicx}	% Inclusão de gráficos
\usepackage{epsfig,subfig}		% Inclusão de figuras
\usepackage{microtype} 			% Melhorias de justificação
% ---------------------
		
% ---------------------
% Pacotes ADICIONAIS
% ---------------------
\usepackage{lipsum}						% Geração de dummy text
\usepackage{amsmath,amssymb,mathrsfs}	% Comandos matemáticos avançados 
\usepackage{setspace}  					% Para permitir espaçamento simples, 1 1/2 e duplo
\usepackage{verbatim}					% Para poder usar o ambiente "comment"
\usepackage{tabularx} 					% Para poder ter tabelas com colunas de largura auto-ajustável
\usepackage{afterpage} 					% Para executar um comando depois do fim da página corrente
\usepackage{url} 						% Para formatar URLs (endereços da Web)
% ---------------------

% ---------------------
% Pacotes de CITAÇÕES
% ---------------------
\usepackage[brazilian,hyperpageref]{backref}	% Paginas com as citações na bibl
\usepackage[alf]{abntex2cite}				% Citações padrão ABNT (alfa)
%\usepackage[num]{abntex2cite}				% Citações padrão ABNT (numericas)
% ---------------------

% Configurações de CITAÇÕES para abntex2
% --- 
% CONFIGURAÇÕES DE PACOTES
% --- 

% ---
% Configurações do pacote backref
% Usado sem a opção hyperpageref de backref
\renewcommand{\backrefpagesname}{Citado na(s) página(s):~}
% Texto padrão antes do número das páginas
\renewcommand{\backref}{}
% Define os textos da citação
\renewcommand*{\backrefalt}[4]{
	\ifcase #1 %
		Nenhuma citação no texto.%
	\or
		Citado na página #2.%
	\else
		Citado #1 vezes nas páginas #2.%
	\fi}%
% ---

% Inclusão de dados para CAPA e FOLHA DE ROSTO (título, autor, orientador, etc.)
% ---
% Informações de dados para CAPA e FOLHA DE ROSTO
% ---
\titulo{Análise Comparativa entre SAST Tradicional e Ferramentas Empresariais de Desenvolvimento Assistido por IA para Detecção de Vulnerabilidades em Aplicações Python}
\autor{Pedro Victor Marcelino Jordão Motta}
\local{Santo André - SP}
\data{Dezembro de 2026}
\orientador{Diogo Santana Martins}
\instituicao{%
  Universidade Federal do ABC -- UFABC
  \par
  Centro de Matemática, Computação e Cognição 
  \par
  Programa de Especialização em Tecnologias e Sistemas de Informação (TSI)}
\tipotrabalho{TCC da Especialização em Tecnologias e Sistemas de Informação (TSI)}
% O preambulo deve conter o tipo do trabalho, o objetivo,
% o nome da instituição e a área de concentração
\preambulo{\textbf{Trabalho de Conclusão em Tecnologias e Sistemas de Informação (TSI)} apresentada ao Programa de Pós-Graduação.}
% ---

% Inclui Configurações de aparência do PDF Final
%  Configurações de aparência do PDF final
% NÃO ALTERAR!!!

% alterando o aspecto da cor azul
\definecolor{blue}{RGB}{41,5,195}

% informações do PDF
\makeatletter
\hypersetup{
     	%pagebackref=true,
		pdftitle={\@title}, 
		pdfauthor={\@author},
    		pdfsubject={\imprimirpreambulo},
	    pdfcreator={LaTeX with abnTeX2},
		pdfkeywords={abnt}{latex}{abntex}{abntex2}{trabalho acadêmico}, 
		colorlinks=true,       		% false: boxed links; true: colored links
    		linkcolor=blue,          	% color of internal links
    		citecolor=blue,        		% color of links to bibliography
    		filecolor=magenta,      		% color of file links
		urlcolor=blue,
		bookmarksdepth=4
} 
\makeatother
% --- 

% O tamanho da identação do parágrafo é dado por:
\setlength{\parindent}{1.3cm}

% Controle do espaçamento entre um parágrafo e outro:
\setlength{\parskip}{0.2cm}  % tente também \onelineskip

% ---------------------
% Compila o indice
% ---------------------
\makeindex
% ---------------------

%%%%%%%%%%%%%%%%%%%%%%%%%%%
%%  INICIO DO DOCUMENTO  %%
%%%%%%%%%%%%%%%%%%%%%%%%%%%
\begin{document}

% Retira espaço extra obsoleto entre as frases.
\frenchspacing

% ----------------------------------------------------------
% ELEMENTOS PRÉ-TEXTUAIS (Capa, Resumo, Abstract, etc.)
% ----------------------------------------------------------
\pretextual

% Capa
% ---
% Impressão da Capa
% ---
  \begin{capa}%
    \begin{figure}[h!]%
        \centering%
        \includegraphics[scale=1.2]{figs/logo.png}%
      \end{figure}%
    \center
	\ABNTEXchapterfont\large{Universidade Federal do ABC \\ Centro de Matemática, Computação e Cognição\\ Programa de Pós-Graduação em Ciência da Computação}
	%\vspace{1.5cm}

    \vfill
    \ABNTEXchapterfont\bfseries\LARGE\imprimirtitulo
    \vfill

	%\vfill
	\ABNTEXchapterfont\large\imprimirautor
	\vfill
%
	
    \large\imprimirlocal, \large\imprimirdata

    \vspace*{1cm}
  \end{capa}
% ---

% Folha de rosto (o * indica que haverá a ficha bibliográfica)
\imprimirfolhaderosto*

% Imprimir Ficha Catalografica
% ---
% Ficha Catalográfica
% ---
% Isto é um exemplo de Ficha Catalográfica, ou ``Dados internacionais de
% catalogação-na-publicação''. Você pode utilizar este modelo como referência. 
% Porém, talvez a biblioteca lhe fornece um PDF
% com a ficha catalográfica definitiva após a defesa do trabalho. Quando estiver
% com o documento, salve-o como PDF no diretório do seu projeto e substitua todo
% o conteúdo de implementação deste arquivo pelo comando abaixo:
%
% \begin{fichacatalografica}
%     \includepdf{fig_ficha_catalografica.pdf}
% \end{fichacatalografica}
\begin{fichacatalografica}
	\vspace*{\fill}					% Posição vertical
	\hrule							% Linha horizontal
	\begin{center}					% Minipage Centralizado
	\begin{minipage}[c]{12.5cm}		% Largura
	
	\imprimirautor
	
	\hspace{0.5cm} \imprimirtitulo  / \imprimirautor. --
	\imprimirlocal, \imprimirdata-
	
	\hspace{0.5cm} \pageref{LastPage} p. : il. (algumas color.) ; 30 cm.\\
	
	\hspace{0.5cm} \imprimirorientadorRotulo~\imprimirorientador\\
	
	\hspace{0.5cm}
	\parbox[t]{\textwidth}{\imprimirtipotrabalho~--~\imprimirinstituicao,
	\imprimirdata.}\\
	
	\hspace{0.5cm}
		1. Palavra-chave1.
		2. Palavra-chave2.
		I. Orientador.
		II. Universidade xxx.
		III. Faculdade de xxx.
		IV. Título\\ 			
	
	\hspace{8.75cm} CDU 02:141:005.7\\
	
	\end{minipage}
	\end{center}
	\hrule
\end{fichacatalografica}
% ---

% Inserir Folha de Aprovação
% ---
% Assinaturas
% ---
% Este é um exemplo de folha de aprovação, elemento obrigatório da NBR
% 14724/2011 (seção 4.2.1.3). Você pode utilizar este modelo até a aprovação
% do trabalho. Após isso, substitua todo o conteúdo deste arquivo por uma
% imagem da página assinada pela banca com o comando abaixo:
%
% \includepdf{folhadeaprovacao_final.pdf}
%
\begin{folhadeaprovacao}

  \begin{center}
    {\ABNTEXchapterfont\large\imprimirautor}

    \vspace*{\fill}\vspace*{\fill}
    \begin{center}
      \ABNTEXchapterfont\bfseries\Large\imprimirtitulo
    \end{center}
    \vspace*{\fill}
    
    \hspace{.45\textwidth}
    \begin{minipage}{.5\textwidth}
        \imprimirpreambulo
    \end{minipage}%
    \vspace*{\fill}
   \end{center}
        
 % Isso na versao final do trabalho!!!       
   Trabalho aprovado. \imprimirlocal, 01 de janeiro de 20XX:

   \assinatura{\textbf{\imprimirorientador} \\ Orientador} 
   \assinatura{\textbf{Professor} \\ Convidado 1}
   \assinatura{\textbf{Professor} \\ Convidado 2}
   \assinatura{\textbf{Professor} \\ Convidado 3}
      
   \begin{center}
    \vspace*{0.5cm}
    {\large\imprimirlocal}
    \par
    {\large\imprimirdata}
    \vspace*{1cm}
  \end{center}
  
\end{folhadeaprovacao}
% ---

% Dedicatória
% ---
% Dedicatória
% ---
\begin{dedicatoria}
   \vspace*{\fill}
   \centering
   \noindent
   \textit{ Aos verme que roeu as frias carnes de meu cadáver.} \vspace*{\fill}
\end{dedicatoria}
% ---

% Agradecimentos
% ---
% Agradecimentos
% ---
\begin{agradecimentos}


Agradeço ao meu orientador, XXXXXXXXX, por todos os conselhos, pela paciência e ajuda nesse período.

Aos meus amigos ...

Aos professores ...

À XXXXXX pelo apoio financeiro para realização deste trabalho de pesquisa.

\end{agradecimentos}
%% ---

% Epígrafe
% ---
% Epígrafe
% ---
\begin{epigrafe}
    \vspace*{\fill}
	\begin{flushright}
		\textit{``Não sei o que, \\
		          não sei o que,\\
                  não sei o que lá.''\\
		          (Autor Desconhecido)}
	\end{flushright}
\end{epigrafe}
% ---

% Resumo e Abstract
% ---
% RESUMOS
% ---

% RESUMO em português
\setlength{\absparsep}{18pt} % ajusta o espaçamento dos parágrafos do resumo
\begin{resumo}

Este trabalho propõe uma análise comparativa entre ferramentas tradicionais de Static Application Security Testing (SAST) e ferramentas empresariais de desenvolvimento assistido por Inteligência Artificial para detecção de vulnerabilidades em aplicações Python. A pesquisa considera três abordagens: SAST tradicional puro, IA pura e uma abordagem híbrida SAST + IA. Serão avaliadas ferramentas como Bandit, Semgrep Community Edition, Cursor e OpenAI Codex, utilizando benchmarks e repositórios vulneráveis em Python. A metodologia envolve revisão bibliográfica, definição de protocolo experimental, execução das ferramentas e comparação dos resultados obtidos. A avaliação será feita por métricas quantitativas, como precisão, recall, F1-score e falsos positivos, além de critérios qualitativos, como clareza da explicação, severidade e utilidade da recomendação. Espera-se identificar em quais cenários a IA pode substituir, complementar ou enriquecer o uso de ferramentas SAST tradicionais em processos de segurança de software.

 \textbf{Palavras-chaves}: Static Application Security Testing; SAST; Inteligência Artificial; segurança de aplicações; vulnerabilidades em Python; agentes de codificação; DevSecOps; Bandit; Semgrep; OpenAI Codex; Cursor.
\end{resumo}

% ABSTRACT in english
\begin{resumo}[Abstract]
 \begin{otherlanguage*}{english}
This work proposes a comparative analysis between traditional Static Application Security Testing (SAST) tools and enterprise AI-assisted development tools for vulnerability detection in Python applications. The research considers three approaches: traditional SAST, pure AI-based analysis, and a hybrid SAST + AI approach. Tools such as Bandit, Semgrep Community Edition, Cursor, and OpenAI Codex will be evaluated using vulnerable Python benchmarks or repositories. The methodology includes a literature review, the definition of a standardized experimental protocol, tool execution, and comparison of the obtained results. The evaluation will consider quantitative metrics, such as precision, recall, F1-score, and false positives, as well as qualitative criteria, including explanation clarity, severity classification, and usefulness of recommendations. The study aims to identify in which scenarios AI can replace, complement, or enrich traditional SAST tools in software security processes.
   \vspace{\onelineskip}
 
   \noindent 
   \textbf{Keywords}: Static Application Security Testing; SAST; Artificial Intelligence; application security; Python vulnerabilities; coding agents; DevSecOps; Bandit; Semgrep; OpenAI Codex; Cursor.
 \end{otherlanguage*}
\end{resumo}

% Lista de ilustrações
\pdfbookmark[0]{\listfigurename}{lof}
\listoffigures*
\cleardoublepage

% Lista de tabelas
\pdfbookmark[0]{\listtablename}{lot}
\listoftables*
\cleardoublepage

% Lista de abreviaturas e siglas
\begin{siglas}
\item[API] \textit{Application Programming Interface}: interface que permite a comunicação entre sistemas, aplicações ou serviços.
  
  \item[CE] \textit{Community Edition}: versão comunitária de uma ferramenta, geralmente gratuita ou aberta.
  
  \item[CI/CD] \textit{Continuous Integration / Continuous Delivery}: conjunto de práticas utilizadas para integração, testes, entrega e implantação contínua de software.
  
  \item[CLI] \textit{Command Line Interface}: interface de linha de comando utilizada para interagir com programas por meio de comandos textuais.
  
  \item[CSV] \textit{Comma-Separated Values}: formato de arquivo textual utilizado para representar dados em estrutura tabular.
  
  \item[CVE] \textit{Common Vulnerabilities and Exposures}: catálogo público de vulnerabilidades conhecidas em softwares, sistemas e bibliotecas.
  
  \item[CWE] \textit{Common Weakness Enumeration}: catálogo de fraquezas comuns de software utilizado para classificar tipos de vulnerabilidades.
  
  \item[DevSecOps] \textit{Development, Security and Operations}: abordagem que integra práticas de desenvolvimento, segurança e operações ao longo do ciclo de vida do software.
  
  \item[F1-score] \textit{F1 Measure}: métrica que combina precisão e revocação em uma única medida de avaliação.
  
  \item[F3-score] \textit{F3 Measure}: variação da métrica F-score que atribui maior peso à revocação, útil em contextos nos quais a não identificação de vulnerabilidades é mais crítica.
  
  \item[FN] \textit{False Negative}: falso negativo; ocorre quando uma vulnerabilidade real não é identificada pela ferramenta.
  
  \item[FP] \textit{False Positive}: falso positivo; ocorre quando a ferramenta indica uma vulnerabilidade inexistente ou não aplicável.
  
  \item[IA] Inteligência Artificial: área da computação voltada ao desenvolvimento de sistemas capazes de executar tarefas associadas à inteligência humana.
  
  \item[JSON] \textit{JavaScript Object Notation}: formato leve de representação de dados, utilizado para armazenar e trocar informações estruturadas.
  
  \item[LLM] \textit{Large Language Model}: modelo de linguagem de grande escala capaz de compreender e gerar texto, código e outros conteúdos.
  
  \item[OWASP] \textit{Open Worldwide Application Security Project}: organização voltada à promoção de boas práticas, materiais e projetos relacionados à segurança de aplicações.
  
  \item[PoC] \textit{Proof of Concept}: prova de conceito; implementação experimental utilizada para validar a viabilidade técnica de uma ideia ou abordagem.
  
  \item[PR] \textit{Pull Request}: solicitação de alteração em um repositório de código, geralmente utilizada para revisão antes da integração ao projeto principal.
  
  \item[RAG] \textit{Retrieval-Augmented Generation}: abordagem em que um modelo de linguagem utiliza informações recuperadas de fontes externas para gerar respostas mais contextualizadas.
  
  \item[SAST] \textit{Static Application Security Testing}: técnica de análise estática de segurança utilizada para identificar vulnerabilidades no código-fonte sem executar a aplicação.
  
  \item[TN] \textit{True Negative}: verdadeiro negativo; ocorre quando a ferramenta corretamente não indica vulnerabilidade onde ela não existe.
  
  \item[TP] \textit{True Positive}: verdadeiro positivo; ocorre quando a ferramenta identifica corretamente uma vulnerabilidade existente.
\end{siglas}



% Inserir o SUMÁRIO
\pdfbookmark[0]{\contentsname}{toc}
\tableofcontents*
\cleardoublepage

% ----------------------------------------------------------
% ELEMENTOS TEXTUAIS (Capítulos)
% ----------------------------------------------------------/in
\textual
% Elementos textuais com numeração arábica
\pagenumbering{arabic}
% Reinicia a contagem do número de páginas
\setcounter{page}{1}

% Inclui cada capitulo da Dissertação
% ----------------------------------------------------------
% Introdução 
% Capítulo sem numeração, mas presente no Sumário
% ----------------------------------------------------------

\chapter[Introdução]{Introdução}

O desenvolvimento de software tem incorporado, de forma crescente, práticas de automação voltadas à produtividade, qualidade e segurança. Nesse contexto, ferramentas de Static Application Security Testing (SAST) são utilizadas para analisar código-fonte antes da execução da aplicação, identificando padrões inseguros, vulnerabilidades conhecidas e possíveis riscos de segurança em fases iniciais do ciclo de desenvolvimento \cite{owasp_sast,semgrep_ce}.

Em projetos corporativos, esse tipo de verificação costuma ser integrado a pipelines de CI/CD, revisões de pull request e políticas de governança de código. Apesar de sua relevância, ferramentas SAST tradicionais apresentam limitações conhecidas, especialmente em relação à quantidade de falsos positivos, dificuldade de compreender contexto de negócio e limitaçã ona interpretação de fluxos distribuídos entre múltiplos arquivos e dependência de regras previamente definidas.

Paralelamente, ferramentas de desenvolvimento assistido por IA passaram de simples autocompletes para agentes de codificação capazes de ler repositórios, editar arquivos, executar comandos, propor correções, revisar código e apoiar fluxos mais complexos de engenharia de software. Esses agentes foram criados em resposta à evolução dos modelos de linguagem aplicados a código e à necessidade de apoiar tarefas cada vez mais complexas, como refatorações, correção de bugs, geração de testes e revisão de alterações. A documentação do Codex descreve a ferramenta como um agente de codificação capaz de ler, editar e executar código, inclusive em ambientes locais, IDEs, aplicações desktop e nuvem \cite{openai_codex_cloud,openai_codex_cli}. O Cursor, por sua vez, posiciona-se como um ambiente de desenvolvimento com compreensão do codebase e recursos de agente para auxiliar o desenvolvedor em tarefas de programação \cite{cursor_home,cursor_agents}.

O interesse pelo tema também se justifica pelo avanço recente da adoção dessas ferramentas na prática profissional. A pesquisa Stack Overflow Developer Survey 2025 indicou que 84\% dos respondentes usam ou planejam usar ferramentas de IA no processo de desenvolvimento, e que 51\% dos desenvolvedores profissionais usam essas ferramentas diariamente \cite{stackoverflow2025}. A OpenAI reportou, em 2026, mais de 2 milhões de desenvolvedores usando o Codex semanalmente e crescimento de 6 vezes no uso do Codex em ambientes ChatGPT Business e Enterprise desde janeiro \cite{openai_codex_pricing_2026}. Além disso, a própria página do Cursor afirma que a ferramenta é utilizada por mais de metade das empresas Fortune 500, indicando relevância empresarial para esse tipo de solução \cite{cursor_home}.

Esse cenário cria uma questão relevante para a área de Segurança da Informação: ferramentas empresariais de desenvolvimento assistido por IA podem detectar vulnerabilidades em aplicações Python com qualidade comparável ou complementar às ferramentas SAST tradicionais? Além disso, uma abordagem híbrida, em que a IA analisa e enriquece os relatórios produzidos por SAST, pode melhorar a utilidade dos achados para desenvolvedores e equipes de segurança?



\section{Motivação}\label{sec:motivacao}


A motivação para o desenvolvimento desta pesquisa está no crescimento do uso de IA no desenvolvimento de software e na necessidade de compreender, de forma empírica, seus impactos na análise de segurança de código. Embora ferramentas de IA estejam sendo cada vez mais utilizadas por desenvolvedores, ainda há incertezas sobre sua confiabilidade, reprodutibilidade e capacidade de atuar em tarefas críticas, como identificação e explicação de vulnerabilidades.

Os potenciais beneficiários da pesquisa são desenvolvedores, equipes de segurança de aplicações, times de DevSecOps, pesquisadores em engenharia de software e organizações que adotam ou pretendem adotar ferramentas de IA no ciclo de desenvolvimento. Desenvolvedores podem se beneficiar de recomendações mais claras e contextualizadas para correção de vulnerabilidades. Equipes de segurança podem se beneficiar de critérios comparativos para decidir quando utilizar SAST tradicional, IA pura ou abordagens híbridas. Organizações podem se beneficiar de evidências para orientar políticas de adoção de IA em processos de segurança de software.

O problema a ser pesquisado pode ser descrito em uma frase: avaliar em que medida ferramentas empresariais de desenvolvimento assistido por IA podem detectar, explicar e priorizar vulnerabilidades em aplicações Python quando comparadas a ferramentas SAST tradicionais e a uma abordagem híbrida SAST + IA.

 Atualmente, a literatura e o mercado tratam o problema principalmente por meio de três abordagens. A primeira é o uso de ferramentas SAST tradicionais, como Bandit e Semgrep Community Edition, baseadas em regras, padrões de código e análise sintática. A segunda é o uso de modelos de linguagem e agentes de codificação para revisão de código e identificação de problemas de segurança. A terceira, ainda em amadurecimento, envolve a combinação entre ferramentas estáticas e IA generativa para apoiar triagem, explicação, classificação e priorização dos achados.

Essas soluções, entretanto, não resolvem completamente o problema. Ferramentas SAST tradicionais tendem a ser reproduzíveis e de baixo custo, mas apodem gerar muitos falsos positivos e ter dificuldade para contextualizar vulnerabilidades. Ferramentas de IA podem produzir explicações mais ricas e recomendações mais próximas da linguagem do desenvolvedor, mas apresentam riscos de alucinação, inconsistência, dependência de prompts e menor previsibilidade. A abordagem híbrida pode combinar pontos fortes das duas alternativas, mas ainda precisa ser avaliada com critérios claros, benchmarks adequados e métricas quantitativas e qualitativas.
\section{Objetivos}\label{sec:objetivos}


\subsection{Objetivo geral} \label{subsec:objetivo-geral}

O objetivo geral deste trabalho é avaliar comparativamente abordagens tradicionais e baseadas em Inteligência Artificial para detecção, explicação e priorização de vulnerabilidades em aplicações Python, considerando ferramentas tradicionais de \textit{Static Application Security Testing} (SAST), ferramentas empresariais de desenvolvimento assistido por IA e uma abordagem híbrida SAST + IA.

\subsection{Objetivos específicos} \label{subsec:objetivos-especificos}

Os objetivos específicos são:

\begin{enumerate}
    \item Definir uma base teórica e técnica sobre SAST, agentes de codificação, segurança de aplicações Python e avaliação de ferramentas de detecção de vulnerabilidades.

    \item Selecionar um conjunto de aplicações, benchmarks ou repositórios Python vulneráveis com informações de referência adequadas para avaliação comparativa.

    \item Projetar um protocolo experimental padronizado para execução das abordagens SAST tradicional, IA pura e SAST + IA.

    \item Construir um pipeline de coleta, normalização e organização dos resultados produzidos pelas ferramentas avaliadas.

    \item Avaliar quantitativamente o desempenho das abordagens por meio de métricas como precisão, \textit{recall}, F1-score, taxa de falsos positivos, quantidade de vulnerabilidades identificadas, tempo de execução e custo estimado.

    \item Avaliar qualitativamente a utilidade dos relatórios gerados, considerando aspectos como clareza das explicações, qualidade das recomendações, indicação de severidade, explorabilidade dos resultados e apoio à tomada de decisão.

    \item Comparar os resultados obtidos e identificar os cenários em que ferramentas baseadas em IA podem substituir, complementar ou enriquecer ferramentas SAST tradicionais.
\end{enumerate}


% PARTE - Define a divisão do documento em partes (Não é obrigatório)
\chapter{Fundamentação Teórica}\label{cap:fundTeoric}
\addcontentsline{toc}{chapter}{Fundamentação Teórica}
Este capítulo estabelece a base conceitual sobre a qual o experimento deste trabalho é construído. 

A Seção ~\ref{sec:seguranca-aplicacoes} trata da segurança de aplicações e da análise estática: parte dos conceitos gerais de vulnerabilidade e das taxonomias CWE/CVE, avança para o funcionamento, as vantagens e as limitações estruturais do \textit{Static Application Security Testing} (SAST) e conclui com a descrição das duas ferramentas SAST  (Bandit e Semgrep Community Edition). 

A Seção ~\ref{sec:llms} apresenta os Modelos de Linguagem de Grande Escala (LLMs): seus fundamentos arquiteturais, sua aplicação ao código-fonte e o mecanismo de \textit{Retrieval-Augmented Generation} (RAG), relevante para entender como ferramentas de codificação recuperam contexto de um repositório. 

A Seção ~\ref{sec:agentes-codificacao} discute os agentes de codificação assistidos por IA - a evolução do paradigma de autocomplete para agentes autônomos, o que caracteriza tecnicamente um agente, as ferramentas empresariais avaliadas (Cursor e OpenAI Codex) e o uso específico de IA para segurança de código, incluindo abordagens híbridas com SAST. 

Por fim, a Seção ~\ref{sec:metricas-avaliacao} define as métricas quantitativas e qualitativas empregadas para avaliar ferramentas de detecção de vulnerabilidades, encerrando com uma contextualização dos benchmarks utilizados na literatura recente, detalhados no Capítulo ~\ref{cap:estadoarte}.


\section{Segurança de Aplicações e Análise Estática}\label{sec:seguranca-aplicacoes}
\subsection{Vulnerabilidades de software: conceitos gerais}\label{subsec:vulnerabilidades-software}
Uma vulnerabilidade de software é uma falha na especificação, no projeto ou na implementação de um sistema que pode ser explorada para comprometer sua confidencialidade, integridade ou disponibilidade \cite{MITRECWE2024}. Essa definição situa o problema em dois planos distintos, msas complementares: o plano da \textit{fraqueza} (weakness) , a característica estrutural do código que torna a exploração possível, por exemplo a ausência de sanitização de uma entrada externa antes de seu uso em uma consulta SQL  e o plano da \textit{vulnerabilidade concreta} (vulnerability), a instância real dessa fraqueza em um produto de software específico, potencialmente já documentada e associada a um identificador público. É justamente essa distinção que estrutura as duas taxonomias descritas na Seção  ~\ref{subsec:cwe-cve}.

A motivação para detectar vulnerabilidades o mais cedo possível no ciclo de desenvolvimento é praticamente consensual na literatura, tanto pelo custo crescente de correção em estágios avançados quanto pelo risco de exploração em produção \cite{Dalaq2025SASTSLR, Shiri2024MLSLR}. Esse princípio conhecido como \textit{shift-left}  motivou o movimento \textbf{DevSecOps}, que propõe integrar atividades de segurança (incluindo análise estática automatizada) diretamente aos pipelines de integração e entrega contínuas (CI/CD), em vez de tratá-las como uma fase de auditoria isolada ao final do desenvolvimento \cite{Myrbakken2017DevSecOps}. A adoção de ferramentas de análise automatizada nas fases iniciais do ciclo de vida do software é, portanto, uma consequência direta dessa mudança de postura \cite{Wadhams2024SASTBarriers}.

O problema ganhou uma nova dimensão com a popularização de assistentes de programação baseados em IA generativa. Um estudo pioneiro sobre o GitHub Copilot verificou que aproximadamente 40\% dos programas gerados em cenários alinhados ao \textit{MITRE Top 25 CWE} continham vulnerabilidades \cite{Pearce2022CopilotSecurity}. De forma complementar, um experimento controlado com usuários demonstrou que desenvolvedores assistidos por IA produziram código menos seguro do que um grupo de controle e, paradoxalmente, relataram maior confiança na segurança do próprio resultado \cite{Perry2023InsecureAICode}. Esses achados indicam uma tendência de crescimento do volume de código potencialmente vulnerável em produção, o que reforça a relevância de mecanismos automatizados tradicionais ou baseados em IA capazes de detectar essas falhas antes da implantação, problema central deste trabalho.

\subsection{CWE (Common Weakness Enumeration) e CVE (Common Vulnerabilities and Exposures)}\label{subsec:cwe-cve}
Para viabilizar a comunicação, a comparação de ferramentas e a construção de conjuntos de dados rotulados, a comunidade de segurança de software adota duas taxonomias mantidas pela MITRE Corporation, complementares entre si. A \textbf{Common Weakness Enumeration (CWE)} é um dicionário de classes abstratas de fraquezas de software, organizado hierarquicamente: cada entrada CWE descreve um \textit{tipo} de falha , por exemplo, CWE-89 (\textit{SQL Injection}), CWE-22 (\textit{Path Traversal}) ou CWE-78 (\textit{OS Command Injection}) independentemente de qualquer produto específico em que ela ocorra \cite{MITRECWE2024}. Já o registro \textbf{Common Vulnerabilities and Exposures (CVE)} cataloga ocorrências concretas de vulnerabilidades, cada uma identificada por um código único (por exemplo, CVE-2024-XXXXX) e tipicamente associada a um ou mais identificadores CWE que descrevem sua natureza \cite{MITRECVE2024}. Em outras palavras, a CWE responde "que tipo de falha é essa?", enquanto a CVE responde "onde e quando essa falha específica ocorreu?".

Essas taxonomias cumprem três papéis neste trabalho. Primeiro, servem de vocabulário comum para categorizar e comparar os achados das três abordagens avaliadas SAST, ferramentas empresariais de IA e a combinação híbrida , permitindo verificar se a cobertura por categoria de vulnerabilidade é equivalente entre elas, um cuidado metodológico recomendado por estudos que apontam que nenhuma ferramenta domina todas as categorias CWE simultaneamente \cite{Esposito2024SASTComparison, Zhu2024SASTAndroid}. Segundo, orientam a curadoria de amostras rotuladas de código Python vulnerável utilizadas no experimento, uma vez que datasets como o VUDENC são construídos precisamente a partir de correções de commits associadas a tipos CWE específicos. Terceiro, complementam-se com classificações de risco de mais alto nível voltadas a aplicações web, como o \textbf{OWASP Top 10}, adotado como base de diversos benchmarks sintéticos citados na literatura de ferramentas híbridas (Seção ~\ref{subsubsec:hibrido-sast-ia}), entre eles o OWASP Benchmark \cite{OWASP2021Top10}. É importante notar, contudo, que essas taxonomias apresentam desbalanceamento de cobertura entre linguagens: mapeamentos sistemáticos da literatura mostram forte concentração de estudos e datasets rotulados em C/C++ e Java, com Python historicamente sub-representado \cite{Kalouptsoglou2023VulnPrediction, Germano2025LLMSLR} lacuna discutida com mais profundidade no Capítulo ~\ref{cap:estadoarte}.


\subsection{Static Application Security Testing (SAST)}\label{subsec:sast}
\subsubsection{Funcionamento e técnicas}\label{subsubsec:sast-funcionamento}
\textit{Static Application Security Testing} (SAST) designa a família de técnicas e ferramentas que analisam o código-fonte, o bytecode ou representações intermediárias de um programa em busca de padrões associados a vulnerabilidades, sem executar o programa \cite{Dalaq2025SASTSLR}. Essa ausência de execução é a característica definidora do paradigma e o distingue de técnicas dinâmicas (DAST, testes de penetração) e de abordagens híbridas de análise (IAST).

Do ponto de vista técnico, as ferramentas SAST tipicamente constroem uma ou mais representações intermediárias do programa  árvores sintáticas abstratas (\textit{Abstract Syntax Trees}, AST), grafos de fluxo de controle (\textit{Control Flow Graphs}) e grafos de fluxo de dados (\textit{Data Flow Graphs})  e aplicam sobre elas regras ou consultas (\textit{queries}) que codificam padrões considerados inseguros \cite{Li2024SAzSTCC, Dalaq2025SASTSLR}. Duas técnicas centrais sustentam a maior parte das análises relevantes para vulnerabilidades de injeção e de manipulação de dados sensíveis:


\begin{itemize}
\item \textbf{Análise sintática (baseada em padrões).} Ferramentas como o Bandit percorrem a AST do programa procurando por padrões estruturais conhecidos como inseguros ,por exemplo, a chamada de \texttt{subprocess} com \texttt{shell=True}, o uso de \texttt{eval()} sobre entrada não confiável, ou a construção de comandos SQL por concatenação de strings. Essa técnica é rápida e não exige modelagem de fluxo de dados, mas tende a gerar tanto falsos positivos (quando o padrão aparece em contexto seguro) quanto falsos negativos (quando a mesma falha é expressa de forma sintaticamente diferente da regra cadastrada).
\item \textbf{Análise de fluxo de dados e \textit{taint analysis}.} Ferramentas baseadas em consultas mais expressivas, como o Semgrep e o CodeQL, modelam o fluxo de valores "contaminados" (\textit{tainted}) desde uma fonte não confiável (\textit{source} , por exemplo, entrada do usuário, argumentos de linha de comando, leitura de arquivo) até um ponto sensível de uso (\textit{sink} , por exemplo, execução de comando, consulta SQL, desserialização), sinalizando um alerta quando esse fluxo ocorre sem passar por uma função de sanitização reconhecida (\textit{sanitizer}) \cite{Li2024SASTCC, Dalaq2025SASTSLR}. Essa técnica é mais precisa para vulnerabilidades de injeção que atravessam múltiplas linhas ou funções, mas seu custo computacional é maior e sua eficácia depende diretamente da completude do catálogo de \textit{sources}, \textit{sinks} e \textit{sanitizers} configurado para a linguagem-alvo.
\end{itemize}

\subsubsection{Vantagens e limitações}\label{subsubsec:sast-vantagens-limitacoes}

Por não exigir execução, o SAST pode ser aplicado desde as fases mais iniciais do ciclo de desenvolvimento ,inclusive sobre código incompleto ou não compilável  cobrir a totalidade do código analisável em uma única execução e apresentar comportamento determinístico e de baixo custo computacional, características que o tornam adequado à integração em pipelines de CI/CD \cite{Dalaq2025SASTSLR, Tamberg2025LLMHarness}.

Suas limitações, contudo, são estruturais e bem documentadas empiricamente, e não meramente incidentais. A primeira delas é a dependência direta da qualidade e da abrangência do conjunto de regras: uma dissecção aprofundada de ferramentas baseadas em consultas (CodeQL, Semgrep) mostrou que a maior parte dos falsos negativos observados decorre de consultas frágeis ou incompletas, e não de limitações do motor de análise em si \cite{Li2024SASTCC}. A segunda é a dificuldade em capturar vulnerabilidades que dependem de lógica de negócio, de configuração externa ou de fluxo interprocedural profundo, cenário em que a análise puramente sintática ou de fluxo local tende a falhar \cite{Zhu2024SASTAndroid}. A terceira é o custo de triagem imposto pelo volume de falsos positivos: um estudo longitudinal com seis ferramentas aplicadas a 116 repositórios C++ observou que a maioria dos alertas de algumas delas permanecia sem qualquer tratamento ao longo de anos, evidenciando baixa acionabilidade prática \cite{Aloraini2019SASTWarnings}.

Um padrão recorrente nas avaliações empíricas recentes é a combinação de precisão relativamente alta com \textit{recall} baixo , o SAST tende a "acertar quando alerta", mas a deixar passar despercebida a maioria das vulnerabilidades reais existentes no código analisado. Em uma comparação controlada de oito ferramentas sobre o benchmark sintético Juliet, mesmo a melhor ferramenta atingiu \textit{recall} médio de aproximadamente 0,58, restrito apenas às categorias de vulnerabilidade que ela declarava cobrir \cite{Esposito2024SASTComparison}. Esse resultado é relevante porque contraria a percepção difundida de que o principal problema do SAST seriam exclusivamente os falsos positivos: a cobertura limitada (falsos negativos) mostra-se igualmente crítica, e possivelmente mais grave em termos de risco residual \cite{Esposito2024SASTComparison, Li2023SASTJava}.

Essas limitações técnicas se traduzem em barreiras de adoção documentadas na indústria. Uma revisão de literatura cobrindo 89 estudos identificou os falsos positivos como a barreira mais citada à adoção de SAST, seguida da baixa qualidade das saídas, do tempo de configuração inicial e do esforço manual de correção \cite{Wadhams2024SASTBarriers}. Em linha semelhante, um estudo qualitativo com 20 profissionais da indústria constatou que a maioria enfrenta dificuldade em diagnosticar os resultados das ferramentas e demanda informações mais finas , como rastreamento do defeito e interpretação da causa , em vez de apenas métricas agregadas \cite{Li2025IndustrySAST}. Um relato de experiência sobre integração de SAST em fluxos de CI/CD, incluindo projetos Python/Django, também identificou notificações pouco claras como obstáculo prático relevante \cite{Chembakottu2025SASTUsability}. Adicionalmente, um estudo com 815 \textit{commits} que introduziram vulnerabilidades mostrou que a combinação de cinco ferramentas SAST sinalizava 78\% dos casos em nível de função ou arquivo, mas grande parte dos alertas individuais era irrelevante para a vulnerabilidade em questão, exigindo triagem humana significativa \cite{Charoenwet2024SecureReview}. Em conjunto, esses resultados indicam que a utilidade prática dos relatórios como clareza, explicação, priorização  constitui uma dimensão de avaliação tão relevante quanto as métricas clássicas de detecção, o que fundamenta a componente qualitativa deste trabalho .

Por fim, estudos que compararam abordagens baseadas em regras com abordagens baseadas em aprendizado de máquina observam capacidades complementares entre os dois paradigmas, mas alertam que combinações ingênuas a simples união dos resultados de ambos e podem elevar o \textit{recall} ao custo de degradar significativamente a precisão global \cite{Croft2021RuleLearningSAST}. Esse é um alerta metodológico diretamente relevante para o desenho da abordagem híbrida deste trabalho, que não deve assumir a priori que combinar SAST e IA produz ganhos automáticos.


\subsection{Ferramentas SAST avaliadas}\label{subsec:ferramentas-sast}
No ecossistema Python, o presente trabalho seleciona duas ferramentas SAST gratuitas e amplamente adotadas, representativas de estratégias técnicas distintas: uma análise sintática específica para a linguagem (Bandit) e uma análise baseada em correspondência de padrões multi-linguagem (Semgrep Community Edition). Uma revisão sistemática recente sobre avaliação de ferramentas SAST aponta que nenhuma ferramenta domina todas as categorias de vulnerabilidade, recomendando justamente o uso combinado de abordagens complementares \cite{Dalaq2025SASTSLR} , critério que orientou essa escolha.


\subsubsection{Bandit}\label{subsubsec:bandit}
O Bandit é um analisador estático desenvolvido especificamente para Python, que percorre a AST do código-fonte aplicando um conjunto de \textit{plugins} (regras) predefinidos, cada um associado a um identificador de teste e, na maioria dos casos, a uma categoria CWE \cite{Dalaq2025SASTSLR}. Por operar diretamente sobre a sintaxe da linguagem, o Bandit não exige compilação e é tipicamente rápido, mas sua cobertura está limitada às regras nativas do projeto, sem suporte nativo a análise de fluxo de dados interprocedural complexa.

As evidências quantitativas disponíveis sugerem que o Bandit, mesmo em configuração padrão ("de fábrica"), é comparativamente mais efetivo do que outras ferramentas gratuitas para falhas típicas do ecossistema Python. Em um benchmark com 108 vulnerabilidades reais extraídas de pacotes PyPI, o Bandit em configuração padrão obteve sensibilidade de aproximadamente 59,7\%, ainda que acompanhada de especificidade baixa; com regras customizadas construídas por um especialista, esse valor subiu para a faixa de 72–81\%, evidenciando forte dependência do conjunto de regras , o mesmo padrão observado para ferramentas baseadas em consultas na Seção ~\ref{subsubsec:sast-vantagens-limitacoes} \cite{Marquer2026SAGA}. Em um estudo sobre correção de código Python gerado por LLMs, o Bandit detectou entre 28\% e 46\% das amostras vulneráveis, taxa de duas a três vezes superior à do CodeQL no mesmo conjunto de dados, o que reforça sua adequação como detector de referência para o ecossistema Python \cite{Alrashedy2025FDSP}.

\subsubsection{Semgrep Community Edition}\label{subsubsec:semgrep-ce}
O Semgrep é uma ferramenta de análise estática multi-linguagem baseada em correspondência de padrões sintáticos (\textit{pattern matching}) sobre a AST, com suporte a regras customizadas escritas em uma sintaxe declarativa próxima ao próprio código-fonte da linguagem-alvo, o que reduz a barreira de criação de novas regras em comparação a linguagens de consulta mais especializadas \cite{Dalaq2025SASTSLR}. A edição \textit{Community} (gratuita, de código aberto) inclui um conjunto de regras públicas mantidas pela comunidade e pela empresa mantenedora, sem os recursos de análise interprocedural mais avançada disponíveis nas edições comerciais.

O desempenho do Semgrep evidenciado na literatura reforça, mais uma vez, a forte dependência do conjunto de regras utilizado. No mesmo benchmark de 108 vulnerabilidades reais de pacotes PyPI citado acima, o Semgrep em configuração padrão obteve sensibilidade de apenas cerca de 28,6\% \cite{Marquer2026SAGA}. Em Java, um estudo com 170 \textit{commits} associados a CVEs mostrou que a taxa de detecção do Semgrep quase triplicou de 15,3\% para 44,7\% , após engenharia manual de regras específicas para o conjunto avaliado \cite{Bennett2024Semgrep}. Esse resultado, obtido fora do ecossistema Python, é relevante para o TCC como alerta metodológico: o desempenho do Semgrep Community Edition em configuração padrão, que é a configuração adotada neste experimento por representar o que está prontamente disponível a um desenvolvedor sem investimento adicional em engenharia de regras, deve ser interpretado como um limite inferior do potencial da ferramenta.

Outras ferramentas do ecossistema , notadamente o CodeQL e o analisador de fluxo de dados Pysa  foram consideradas, mas não compõem o braço SAST deste experimento. O CodeQL, em uma análise longitudinal sobre CVEs reais (incluindo projetos Python), apresentou taxa de detecção agregada de apenas 4,4\%, além de regressões de detecção entre versões sucessivas da ferramenta \cite{Noirot2025CodeQLLongitudinal}; o Pysa, no benchmark de pacotes PyPI já citado, apresentou sensibilidade inferior a 1\% em configuração padrão \cite{Marquer2026SAGA}. Bandit e Semgrep Community Edition foram, portanto, priorizados por serem gratuitos, amplamente adotados na prática e representativos de duas estratégias técnicas distintas e complementares.


\section{Modelos de Linguagem de Grande Escala (LLMs)}\label{sec:llms}
\subsection{Conceitos fundamentais}\label{subsec:llms-conceitos}
\subsubsection{Arquitetura Transformer}\label{subsubsec:transformer}
Os LLMs relevantes para este trabalho - os modelos por trás de assistentes de programação como GitHub Copilot, e das ferramentas empresariais avaliadas no experimento, Cursor e OpenAI Codex — são construídos sobre a arquitetura \textbf{Transformer}, proposta por \citeonline{Vaswani2017Transformer}. Sem aprofundar a formulação matemática, que foge ao escopo deste trabalho, cabe destacar a característica central da arquitetura: em vez de processar uma sequência de texto (ou de código) de forma recorrente, token a token, o Transformer utiliza um mecanismo de \textbf{atenção} (\textit{self-attention}) que permite a cada elemento da sequência ponderar diretamente sua relação com todos os demais elementos, independentemente da distância entre eles. Essa propriedade viabilizou o treinamento eficiente e paralelizável de modelos com bilhões de parâmetros e é a base arquitetural comum a praticamente todos os LLMs de propósito geral e especializados em código utilizados atualmente, incluindo os modelos que sustentam as ferramentas avaliadas neste TCC.

\subsubsection{Treinamento, tokenização e janela de contexto}\label{subsubsec:tokenizacao-contexto}
O treinamento de um LLM ocorre tipicamente em duas etapas principais: um \textbf{pré-treinamento} auto-supervisionado sobre grandes volumes de texto (e, no caso de modelos de código, também código-fonte), em que o modelo aprende a prever o próximo token de uma sequência, seguido de etapas de ajuste fino (\textit{fine-tuning}) e alinhamento , por exemplo, via aprendizado por reforço com retroalimentação humana , voltadas a tornar o modelo mais útil e seguro em interações reais. Antes de ser processado pelo modelo, o texto de entrada é convertido em uma sequência de \textbf{tokens} que são unidades que podem corresponder a palavras inteiras, subpalavras ou símbolos de código, conforme o vocabulário do tokenizador utilizado e é sobre essa sequência de tokens que o mecanismo de atenção.

A \textbf{janela de contexto} (\textit{context window}) é o número máximo de tokens que o modelo pode processar simultaneamente em uma única interação, somando a instrução do usuário, o histórico da conversa e todo o conteúdo de código fornecido como contexto. Essa limitação é diretamente relevante para a análise de vulnerabilidades em nível de repositório: quanto maior a janela de contexto, maior a quantidade de código-fonte (múltiplos arquivos, dependências, histórico de chamadas de função) que pode ser fornecida ao modelo em uma única análise. A literatura recente, porém, mostra que uma janela de contexto ampla não elimina o problema de utilização eficiente desse contexto , visto que modelos degradam sua capacidade de localizar informações relevantes conforme o comprimento do texto de entrada aumenta, fenômeno descrito na literatura como \textit{lost-in-the-end} no contexto específico de detecção de vulnerabilidades \cite{Sovrano2025LostInEnd}.


\subsection{LLMs aplicados a código}\label{subsec:llms-codigo}
\subsubsection{Pré-treinamento em código-fonte}\label{subsubsec:pre-treinamento-codigo}
A aplicação de LLMs a tarefas de programação foi viabilizada pelo pré-treinamento (ou pré-treinamento contínuo) sobre grandes volumes de código-fonte público, tipicamente extraído de repositórios abertos. O marco histórico dessa linha é o \textbf{Codex}, um GPT ajustado sobre código do GitHub e avaliado quanto à correção funcional de programas gerados a partir de descrições em linguagem natural; o Codex é o modelo que originou o GitHub Copilot e, por extensão, deu nome à atual ferramenta "OpenAI Codex" avaliada neste trabalho \cite{Chen2021Codex}. Na mesma linha histórica, modelos abertos como o \textbf{Code Llama}, derivado do Llama 2 da Meta AI via pré-treinamento contínuo em código, com variantes especializadas em preenchimento de lacunas (\textit{infilling}) e em janelas de contexto estendidas, ilustram a consolidação do paradigma de modelos de propósito específico para código \cite{Roziere2023CodeLlama}. Esses modelos são citados neste trabalho como contexto histórico e técnico, e não como objeto de avaliação empírica direta — o experimento deste TCC avalia as ferramentas empresariais construídas sobre modelos dessa linhagem, não os modelos isoladamente.


\subsubsection{Capacidades e limitações}\label{subsubsec:llms-capacidades-limitacoes}
A literatura documenta capacidades relevantes de LLMs em tarefas relacionadas a código , como geração, explicação, transformação e, de interesse direto para este trabalho, identificação de padrões potencialmente inseguros , acompanhadas de limitações sistemáticas que precisam ser consideradas na interpretação de qualquer avaliação empírica. A principal delas é a \textbf{alucinação}: a geração de afirmações plausíveis, mas factualmente incorretas ou não fundamentadas no código realmente fornecido, fenômeno que no contexto de detecção de vulnerabilidades se manifesta como falsos positivos com justificativas aparentemente coerentes, mas tecnicamente equivocadas. A segunda é a \textbf{dependência de prompt}: pequenas variações na formulação da instrução produzem diferenças substanciais de desempenho. Estudos citados no Capítulo ~\ref{cap:estadoarte} relatam ganhos de até 0,18 em F1-score entre estratégias de \textit{prompt} \cite{Khare2025LLMEffectiveness} e diferenças superiores a 0,3 em F1-score entre formulações de instrução \cite{Iranmanesh2025ZeroFalse} , o que compromete a reprodutibilidade das avaliações se o protocolo experimental não fixar e documentar cuidadosamente os \textit{prompts} utilizados. A terceira é a \textbf{previsibilidade limitada} do comportamento do modelo entre execuções: mesmo sob a mesma entrada, LLMs comerciais apresentam variação de resultado entre chamadas idênticas, um comportamento não-determinístico documentado já nas primeiras comparações entre LLMs e SAST \cite{Ozturk2023AIvsSAST}, discutido com mais profundidade no Capítulo ~\ref{cap:estadoarte}.

\subsection{Retrieval-Augmented Generation (RAG)}\label{subsec:rag}
Um LLM de propósito geral, por si só, tem acesso apenas ao conhecimento codificado em seus parâmetros durante o treinamento, o que inclui uma amostra de código público, mas não o conteúdo específico do repositório privado que se deseja analisar. O paradigma de \textbf{Retrieval-Augmented Generation (RAG)}, formalizado por \citeonline{Lewis2020RAG}, propõe combinar o modelo generativo com um componente de recuperação (\textit{retriever}) que busca, em uma base externa, os trechos mais relevantes para a tarefa em questão, fornecendo-os ao modelo como contexto adicional antes da geração da resposta. O objetivo é reduzir a alucinação e permitir que o modelo responda com base em informação concreta e verificável, em vez de depender exclusivamente do que foi memorizado durante o treinamento.

Esse princípio é diretamente relevante para entender como as ferramentas de codificação avaliadas neste trabalho operam sobre um repositório: em vez de receber todo o código-fonte de uma vez (o que seria inviável para repositórios grandes, dada a limitação de janela de contexto), essas ferramentas implementam mecanismos de busca por indexação semântica de arquivos, busca textual, navegação guiada pela estrutura do projeto para localizar e fornecer ao modelo apenas os trechos de código considerados relevantes para a tarefa corrente. É importante, contudo, não superestimar essa analogia: ferramentas comerciais de codificação não implementam necessariamente RAG na forma canônica descrita por \citeonline{Lewis2020RAG} , mas sim variações e heurísticas proprietárias inspiradas no mesmo princípio geral de fundamentar a geração em contexto recuperado. Essa distinção é relevante porque a forma como o contexto é entregue ao modelo afeta diretamente seu desempenho em tarefas de análise de segurança: avaliações citadas no Capítulo ~\ref{cap:estadoarte} mostram que abordagens que fornecem contexto estruturado de fatias de código guiadas por grafos de propriedade \cite{Lekssays2025LLMxCPG}, abstrações de funções chamadas \cite{Yang2025ContextVD} ou especificações extraídas de histórico de correções \cite{Zhu2025SpecGuidedVD} , melhoram consistentemente os resultados de detecção em relação a análises de trechos de código isolados, o que reforça a relevância de avaliar ferramentas com acesso real ao repositório, e não apenas modelos respondendo sobre um trecho isolado de código.


\section{Agentes de Codificação Assistidos por IA}\label{sec:agentes-codificacao}
\subsection{De autocomplete a agentes: evolução do paradigma}\label{subsec:evolucao-agentes}
O uso de LLMs no desenvolvimento de software evoluiu em estágios sucessivos de autonomia. Na primeira geração, ferramentas como o GitHub Copilot atuavam como \textbf{completadores de código} (\textit{autocomplete}): a partir do contexto imediato (o texto já digitado no arquivo), o modelo sugeria a continuação mais provável de uma linha ou um pequeno bloco , e cada sugestão exigia aceitação explícita do desenvolvedor antes de ser incorporada ao código. Nessa fase, o LLM não tinha iniciativa: reagia a um contexto local sem planejar múltiplos passos nem interagir com o ambiente além da própria sugestão de texto.

Uma segunda geração introduziu assistentes conversacionais integrados ao ambiente de desenvolvimento, capazes de gerar blocos de código mais extensos, ou de propor alterações em múltiplos trechos de um arquivo, a partir de instruções em linguagem natural .Ainda, contudo, sob iniciativa do desenvolvedor a cada interação. O terceiro estágio, no qual se inserem as ferramentas avaliadas neste trabalho, é o de \textbf{agentes de codificação}: sistemas em que o modelo decompõe uma tarefa de alto nível em subtarefas, decide autonomamente quais ferramentas invocar (busca no repositório, leitura e escrita de arquivos, execução de comandos de terminal, execução de testes ou de analisadores), avalia os resultados obtidos e itera até considerar a tarefa concluída , com supervisão humana em pontos de controle definidos, mas sem exigir aprovação a cada micro-ação. A consolidação desse paradigma em engenharia de software foi impulsionada por benchmarks de resolução autônoma de \textit{issues} reais de repositório, como o SWE-bench \cite{Jimenez2024SWEbench}, e pela demonstração de que a qualidade da interface entre o agente e o ambiente computacional impacta fortemente seu desempenho \cite{Yang2024SWEagent}.


\subsection{O que caracteriza um ``agente''}\label{subsec:caracterizacao-agente}
O ciclo de operação de um agente de LLM é formalizado, na literatura, pelo paradigma \textbf{ReAct} (\textit{Reasoning and Acting}): o modelo intercala passos explícitos de raciocínio (\textit{thought})  em que articula, em linguagem natural, o que precisa ser feito e por quê , com ações concretas (\textit{action}) e invocações de ferramentas externas ao próprio modelo , além de observações (\textit{observation}) do resultado dessas ações no ambiente, repetindo esse ciclo até concluir a tarefa ou atingir um critério de parada \cite{Yao2023ReAct}. Três elementos, tomados em conjunto, distinguem um agente de um simples completador de texto ou de um chatbot de pergunta-e-resposta:

1. \textbf{Uso de ferramentas (\textit{tool-calling}).} O modelo pode invocar funções externas, como leitura e escrita de arquivos, execução de comandos de shell, chamadas a analisadores estáticos ou a interpretadores  e incorporar o resultado dessas chamadas ao seu raciocínio subsequente, em vez de responder apenas com base em seu conhecimento paramétrico.

2. \textbf{Iteração autônoma.} O modelo decide, por conta própria, quantos ciclos de raciocínio-ação são necessários para concluir a tarefa, sem que cada passo intermediário exija uma nova instrução explícita do usuário.

3. \textbf{Acesso e leitura do estado do ambiente.} O agente pode explorar a estrutura de um repositório, ler múltiplos arquivos relacionados e navegar entre eles conforme a necessidade identificada durante o próprio raciocínio, em vez de depender exclusivamente do trecho de código fornecido a priori pelo usuário.

A formalização de \textbf{interfaces agente-computador} (\textit{Agent-Computer Interfaces}) , o conjunto de ferramentas e formatos de retorno disponibilizados ao agente demonstrou que a qualidade dessa interface afeta diretamente o desempenho observável: com uma interface especializada, a taxa de resolução de \textit{issues} reais no benchmark SWE-bench passou de 1,31\% (quando o modelo apenas recuperava contexto de forma simples) para 12,47\% \cite{Yang2024SWEagent, Jimenez2024SWEbench}. Ainda assim, as taxas absolutas de resolução permanecem baixas em tarefas gerais de engenharia de software, o que calibra expectativas realistas sobre a autonomia desses sistemas em tarefas complexas de repositório \cite{Jimenez2024SWEbench} , calibração especialmente relevante para a interpretação dos resultados deste TCC, cujo domínio (detecção de vulnerabilidades) é tipicamente mais difícil de verificar automaticamente do que a resolução de um \textit{bug} funcional com um conjunto de testes associado.

\subsection{Modelos, agentes e camada operacional}\label{subsec:camada-operacional-agente}
Para interpretar corretamente o objeto de avaliação deste trabalho, é necessário distinguir três níveis. O \textbf{LLM isolado} recebe uma entrada delimitada e produz texto, uma classificação ou chamadas de ferramenta; por si só, não explora um repositório nem executa ações no ambiente. O \textbf{agente de codificação} organiza ciclos de raciocínio, ação e observação para cumprir uma tarefa, conforme descrito na Seção ~\ref{subsec:caracterizacao-agente}. Por fim, a \textbf{camada operacional do agente} conecta o modelo ao ambiente: seleciona e recupera contexto, mantém o estado da tarefa, disponibiliza ferramentas, executa comandos, aplica permissões e define critérios de parada. A literatura recente também denomina essa camada de \textit{agent harness}, embora o termo ainda não seja plenamente padronizado \cite{Yang2024SWEagent, Macedo2026AgentHarness}.

Cursor e OpenAI Codex não são, portanto, modelos de IA isolados nem devem ser tratados como sinônimo de "IA pura". São ferramentas empresariais de desenvolvimento assistido por IA que implementam agentes de codificação apoiados por uma camada operacional própria. Seus resultados refletem a configuração completa ferramenta-agente onde contém modelo(s), recuperação de contexto, ferramentas disponíveis, permissões e iteração  e não somente uma capacidade intrínseca de um modelo de linguagem. Neste trabalho, a expressão \textbf{LLM isolado} é reservada aos estudos em que o modelo recebe código em uma entrada fixa, sem exploração instrumental autônoma do repositório; Cursor e Codex constituem o braço de \textbf{ferramentas-agente empresariais}.

\subsection{Ferramentas avaliadas}\label{subsec:ferramentas-agentes}
Este trabalho avalia duas ferramentas empresariais de desenvolvimento assistido por IA que implementam o paradigma de agente descrito nas Seções ~\ref{subsec:evolucao-agentes} e ~\ref{subsec:caracterizacao-agente}, operando com acesso direto ao repositório de código sob análise.

\subsubsection{Cursor}\label{subsubsec:cursor}
O Cursor é um editor de código construído sobre a base do Visual Studio Code, que incorpora um modo de operação agêntico (\textit{Agent Mode}) no qual o modelo compreende a estrutura do repositório, propõe e executa alterações em múltiplos arquivos, executa comandos de terminal e itera sobre o resultado dessas ações até considerar a tarefa concluída, mediante instruções em linguagem natural \cite{Cursor2025Docs}. No contexto deste trabalho, o Cursor é operado em modo agente para realizar a análise de segurança do código Python fornecido, com acesso de leitura ao repositório e, conforme o protocolo experimental definido na metodologia, possivelmente também a ferramentas de execução para validação de hipóteses sobre o comportamento do código.

\subsubsection{OpenAI Codex}\label{subsubsec:openai-codex}
O OpenAI Codex, na configuração avaliada neste trabalho, refere-se ao agente de codificação de linha de comando (\textit{Codex CLI}) disponibilizado pela OpenAI, que executa localmente com acesso configurável a arquivos, comandos de terminal e políticas de permissão (\textit{sandbox}) \cite{OpenAI2025CodexDocs}. É importante notar a ambiguidade histórica do nome: "Codex" designou originalmente o modelo de linguagem ajustado para código descrito por \citeonline{Chen2021Codex}, hoje descontinuado como produto isolado; a ferrasmenta atual reutiliza o nome, mas opera como um agente autônomo construído sobre modelos mais recentes, e não como o completador de código original. Essa distinção é registrada aqui para evitar confusão entre o objeto de estudo empírico deste TCC (a ferramenta-agente atual) e a referência histórica citada na Seção ~\ref{subsubsec:pre-treinamento-codigo}.


\subsection{Uso de IA para segurança de código}\label{subsec:ia-seguranca-codigo}
\subsubsection{Detecção de vulnerabilidades via LLM}\label{subsubsec:deteccao-llm}
Para além da geração de código, LLMs vêm sendo empregados diretamente como detectores de vulnerabilidades, seja por meio de \textit{prompts} que solicitam a análise de um trecho de código, seja por meio de agentes com acesso ao repositório. Surveys recentes documentam a rápida expansão dessa linha de pesquisa, catalogando de 58 a 208 estudos primários sobre o tema \cite{Sheng2025LLMSurvey, Zhou2025LLMReview, Germano2025LLMSLR}. O padrão geral observado que será detalhado com maior profundidade no Capítulo ~\ref{cap:estadoarte} , é de potencial de cobertura mais amplo que o das ferramentas SAST tradicionais, acompanhado de taxas de falso positivo mais altas, sensibilidade a variações de \textit{prompt} e de benchmark, e desempenho substancialmente reduzido quando avaliado sob condições metodologicamente mais rigorosas (dados sem vazamento entre treino e teste, comparações em pares entre versão vulnerável e corrigida) \cite{Ding2025CodeLMVuln, Khare2025LLMEffectiveness}.

Estudos de campo indicam, adicionalmente, uma distância relevante entre o desempenho reportado em benchmark e a utilidade percebida por desenvolvedores reais: em um estudo com 17 desenvolvedores utilizando uma ferramenta de detecção e reparo baseada em IA integrada ao ambiente de desenvolvimento, os alertas gerados receberam utilidade média de 2,5 em uma escala de 5 pontos, e 75\% das correções sugeridas não eram diretamente aplicáveis \cite{Steenhoek2025IDEUserStudy}. Esse resultado reforça, no contexto específico de ferramentas de IA, a mesma conclusão já observada para SAST na Seção ~\ref{subsubsec:sast-vantagens-limitacoes}: métricas de detecção favoráveis não implicam, automaticamente, relatórios úteis na prática de desenvolvimento — o que fundamenta a avaliação qualitativa conduzida neste trabalho (Seção ~\ref{subsec:metricas-qualitativas}).

\subsubsection{Abordagens híbridas SAST + IA}\label{subsubsec:hibrido-sast-ia}
Diante do padrão de \textit{trade-off} entre SAST (preciso, porém de cobertura limitada) e LLMs (cobertura potencialmente maior, porém com mais falsos positivos e menor determinismo), uma linha de pesquisa crescente propõe combinar as duas abordagens. O padrão arquitetural mais próximo do desenho experimental deste TCC posiciona o LLM a jusante do SAST: a ferramenta estática gera os alertas iniciais, e o LLM os revisa, filtra, explica ou valida com base no contexto do código, atuando como uma camada de triagem inteligente. Resultados reportados nessa linha incluem reduções substanciais de falsos positivos mantendo a maior parte dos verdadeiros positivos \cite{Ameen2026QASecClaw, Iranmanesh2025ZeroFalse, Agrawal2025SASTGenius}, ainda que nem sempre de forma uniforme entre categorias de vulnerabilidade — em alguns cenários, a filtragem por LLM chega a piorar o desempenho em categorias específicas, e o SAST puro permanece preferível quando o \textit{recall} é prioridade absoluta \cite{Ciavatta2026LLMSAST}.

Outros padrões arquiteturais posicionam o LLM a montante do motor estático , por exemplo, inferindo especificações de \textit{taint} (pares \textit{source}/\textit{sink}) para alimentar ferramentas como o CodeQL \cite{Li2025IRIS} . Ou utilizam o SAST como oráculo de realimentação em laços iterativos de geração e correção de código, abordagem já demonstrada especificamente em Python \cite{Alrashedy2025FDSP, Lee2026CodeEnhancer}. Esses três padrões de filtragem a jusante, ampliação a montante e realimentação em laço  são discutidos com maior detalhamento empírico, incluindo as ressalvas metodológicas sobre a maturidade ainda incipiente dessa literatura (concentração em preprints, escassa validação em Python), no Capítulo ~\ref{cap:estadoarte}.


\section{Métricas de Avaliação de Ferramentas de Detecção}\label{sec:metricas-avaliacao}
\subsection{Métricas quantitativas}\label{subsec:metricas-quantitativas}
A avaliação quantitativa de detectores de vulnerabilidades , sejam eles baseados em regras (SAST) ou em LLMs, utiliza majoritariamente métricas derivadas da matriz de confusão: \textbf{precisão} (proporção de alertas emitidos que correspondem a vulnerabilidades reais), \textbf{recall} ou sensibilidade (proporção de vulnerabilidades reais efetivamente detectadas), \textbf{F1-score} (média harmônica entre precisão e recall) e a \textbf{taxa de falsos positivos} \cite{Kalouptsoglou2023VulnPrediction, Esposito2024SASTComparison}. Um mapeamento sistemático com 180 estudos identificou o F1-score como a métrica mais popular, mas recomendou o uso complementar de variantes que penalizam mais fortemente os falsos negativos, como o \textbf{F2-score}, sob o argumento de que, em segurança, deixar de detectar uma vulnerabilidade real tende a ser mais grave do que gerar um alerta espúrio que será descartado na triagem \cite{Kalouptsoglou2023VulnPrediction}. Na mesma direção, benchmarks recentes adotam o \textbf{F3-score} como métrica primária, penalizando ainda mais os falsos negativos \cite{Pellew2026RealVuln}, e a revisão sistemática de \citeonline{Dalaq2025SASTSLR} recomenda o \textbf{índice de Youden} como medida balanceada entre sensibilidade e especificidade, menos sensível ao desbalanceamento de classes do que o F1 isolado.

A literatura recente alerta, contudo, que métricas agregadas podem ser enganosas quando calculadas sobre benchmarks metodologicamente frágeis. Foi demonstrado que rótulos incorretos e vazamento de dados entre conjuntos de treino e teste inflam artificialmente os resultados reportados: um modelo com F1 de 68,3\% em um benchmark tradicional caiu para 3,1\% quando reavaliado em um conjunto de dados corrigido e deduplicado \cite{Ding2025CodeLMVuln}. Como resposta a esse problema, foram propostas métricas mais rigorosas, como o \textbf{VD-S} (taxa de falsos negativos sob uma taxa de falsos positivos limitada) e a \textbf{avaliação em pares} (comparar o julgamento do detector sobre a função vulnerável e sobre sua versão corrigida), que penaliza detectores que simplesmente classificam tudo como vulnerável, prática que infla artificialmente o \textit{recall} aparente \cite{Ding2025CodeLMVuln, Yildiz2025LLMAgents}.


\subsection{Métricas qualitativas}\label{subsec:metricas-qualitativas}
Métricas puramente quantitativas não capturam integralmente a utilidade de um relatório de vulnerabilidade para um desenvolvedor em contexto real de trabalho, como evidenciado repetidamente ao longo deste capítulo (Seções ~\ref{subsubsec:sast-vantagens-limitacoes} e ~\ref{subsubsec:deteccao-llm}). Cresce, por isso, o reconhecimento de que a utilidade prática dos relatórios deve ser avaliada explicitamente, por meio de dimensões qualitativas complementares:

\begin{itemize}
\item \textbf{Clareza da explicação} — se o relatório comunica, em linguagem compreensível, por que o trecho de código é considerado vulnerável, e não apenas qual regra ou padrão foi acionado.
\item \textbf{Qualidade da recomendação de correção} — se a sugestão fornecida (quando existente) é tecnicamente correta, específica ao contexto do código e diretamente aplicável, ou genérica a ponto de exigir retrabalho substancial do desenvolvedor.
\item \textbf{Indicação de severidade e priorização} — se o relatório ajuda a hierarquizar múltiplos achados por criticidade, facilitando a triagem em cenários com grande volume de alertas, precisamente o gargalo identificado como principal barreira de adoção do SAST na Seção ~\ref{subsubsec:sast-vantagens-limitacoes}.
\end{itemize}
Rubricas qualitativas formais foram propostas para operacionalizar essa avaliação, notadamente a rubrica de \textbf{corretude, clareza e acionabilidade} (\textit{accuracy, clarity, actionability}) \cite{Mao2025ExplainableVD}, que constatou que uma alta taxa de detecção não implica, por si só, explicações acionáveis. Estudos com profissionais da indústria corroboram essa dissociação: alertas tecnicamente corretos podem ser percebidos como pouco úteis no fluxo real de trabalho \cite{Steenhoek2025IDEUserStudy, Li2025IndustrySAST}. Este trabalho combina, portanto, métricas quantitativas (precisão, \textit{recall}, F1, taxa de falsos positivos) com avalfiação qualitativa dos relatórios produzidos pelos três braços do experimento (clareza, qualidade da recomendação e indicação de severidade), seguindo esse consenso metodológico emergente na literatura.

\subsection{Benchmarks de vulnerabilidades}\label{subsec:benchmarks-vulnerabilidades}
A validade de qualquer comparação quantitativa entre ferramentas depende diretamente da qualidade do benchmark utilizado , isto é, do conjunto de código rotulado como vulnerável ou seguro sobre o qual as métricas da Seção ~\ref{subsubsec:bandit} são calculadas. A literatura recente demonstra que benchmarks amplamente utilizados sofrem de problemas metodológicos sistemáticos: rótulos incorretos, vazamento de dados entre partições de treino e teste, e distribuições de vulnerabilidade pouco realistas em relação a código de produção, todos fatores que inflam artificialmente os resultados reportados \cite{Ding2025CodeLMVuln, Chen2023DiverseVul, Chakraborty2022DLLimits}. Avaliações conduzidas em nível de função isolada tendem, adicionalmente, a superestimar a capacidade real dos detectores em comparação a avaliações em nível de repositório completo, onde o contexto necessário para julgar corretamente um trecho de código está frequentemente disperso em múltiplos arquivos \cite{Wen2026VulEval}.

Diante desse cenário, benchmarks mais recentes e cuidadosamente curados assumem papel central na avaliação empírica deste trabalho, com destaque para o \textbf{RealVuln}, benchmark focado especificamente no ecossistema Python, com dados rotulados a partir de vulnerabilidades reais em repositórios de código aberto \cite{Pellew2026RealVuln}. Sua metodologia de construção, seus resultados de comparação entre ferramentas de regras, LLMs de propósito geral e \textit{scanners} especializados, e sua relação direta com o protocolo experimental adotado neste TCC são detalhados no Capítulo ~\ref{cap:estadoarte}, no contexto mais amplo do estado da arte em comparações entre SAST e abordagens baseadas em IA.

\chapter{Estado da Arte}\label{cap:estadoarte}
\addcontentsline{toc}{chapter}{Estado da Arte}

Este capítulo consolida os trabalhos recentes diretamente relacionados ao problema deste TCC: uma comparação empírica, sob protocolo único e em aplicações Python, entre SAST tradicional, ferramentas empresariais de desenvolvimento assistido por IA e uma abordagem híbrida SAST + IA. Cada seção isola uma peça desse desenho experimental e termina apontando o que falta para viabilizá-lo. 

A Seção ~\ref{sec:avaliacoes-sast} mostra como ferramentas SAST têm sido avaliadas e por que seu comportamento real é o ponto de partida (e o baseline) do experimento. 

A Seção ~\ref{sec:llms-deteccao-vulnerabilidades} mostra como LLMs isolados vêm sendo aplicados à mesma tarefa, distinguindo-os das ferramentas-agente avaliadas neste TCC. 

A Seção ~\ref{sec:comparacoes-sast-ia} revisa os estudos que já colocam SAST e LLMs lado a lado , a linha de pesquisa mais próxima deste trabalho  e expõe o quanto ela ainda ignora Python e ferramentas comerciais. 

A Seção ~\ref{sec:abordagens-hibridas}  trata do terceiro braço, o híbrido, mapeando os padrões de combinação já testados. 
s
A Seção ~\ref{sec:agentes-ia-seguranca} examina o paradigma de agentes de IA em segurança de código, do qual emergem Cursor e OpenAI Codex; Claude Code é tratado apenas como alternativa comercial contextual. 

A Seção ~\ref{sec:lacunas-literatura} sintetiza as lacunas remanescentes — inclusive o resultado de uma busca complementar dedicada, especificamente para verificar se a lacuna mais crítica identificada neste levantamento já havia sido preenchida pela literatura mais recente.


\section{Avaliações recentes de ferramentas SAST}\label{sec:avaliacoes-sast}

Antes de comparar SAST com qualquer abordagem baseada em IA, é preciso estabelecer o que o SAST tradicional efetivamente entrega quando avaliado com rigor , esse é o baseline contra o qual os braços de ferramentas-agente empresariais e híbrido deste TCC serão medidos, e por isso a literatura revisada nesta seção é fundação direta do desenho experimental. As avaliações empíricas recentes convergem em um diagnóstico: ferramentas SAST tendem a ser precisas quando alertam, mas deixam de detectar a maior parte das vulnerabilidades reais. Em benchmark sintético de larga escala, a melhor entre oito ferramentas atingiu \textit{recall} médio de aproximadamente 0,58 mesmo nas categorias que declarava cobrir \cite{Esposito2024SASTComparison}. Em código real, o quadro é mais severo: sete ferramentas gratuitas para Java, que atingiam F1 de até 84\% em benchmark sintético, detectaram no máximo 12,7\% de 165 CVEs reais, e mesmo a combinação de quatro ferramentas deixou 70,9\% dos casos sem detecção \cite{Li2023SASTJava}. Avaliação de quatro ferramentas em 170 \textit{commits} com CVEs confirmou taxas individuais entre 11\% e 27\%, com a combinação alcançando 38,8\% \cite{Bennett2024Semgrep}, e um estudo com nove ferramentas para JavaScript observou taxa média de verdadeiros positivos de 15,1\%, com precisão extremamente baixa nas ferramentas de maior cobertura \cite{Brito2023JSSAST}. No maior estudo longitudinal disponível, o CodeQL detectou apenas 4,4\% de quase 4 mil CVEs analisáveis, incluindo apenas 36 casos em Python, além de apresentar regressões de detecção entre versões \cite{Noirot2025CodeQLLongitudinal}.

Dois fatores explicativos se repetem. O primeiro é a dependência do conjunto de regras: a taxa de detecção do Semgrep quase triplicou (15,3\% para 44,7\%) após engenharia manual de regras \cite{Bennett2024Semgrep}, e a dissecção de falsos negativos em ferramentas baseadas em consultas mostrou que a maioria decorre de capacidade inadequada das regras, e não do motor de análise \cite{Li2024SASTCC, Li2023SASTJava}. Em Python, o mesmo padrão foi observado: Bandit e Semgrep em configuração padrão detectaram, respectivamente, cerca de 59,7\% e 28,6\% de 108 vulnerabilidades reais de pacotes PyPI, valores que subiram para 72–81\% com regras customizadas construídas ao custo de meses de trabalho especializado \cite{Marquer2026SAGA}. O segundo fator é a complementaridade entre ferramentas: nenhuma domina todas as categorias, e a combinação eleva a cobertura de forma consistente \cite{Esposito2024SASTComparison, Zhu2024SASTAndroid, Dalaq2025SASTSLR}.

Paralelamente, a linha de pesquisa sobre adoção reforça que o custo de triagem é o gargalo prático: falsos positivos são a barreira mais citada na literatura \cite{Wadhams2024SASTBarriers}, grande parte dos alertas permanece anos sem tratamento \cite{Aloraini2019SASTWarnings}, e profissionais da indústria demandam diagnósticos mais finos do que métricas agregadas \cite{Li2025IndustrySAST}. Em conjunto, esses estudos estabelecem o perfil do braço SAST do presente experimento determinístico, rápido e barato, porém com cobertura limitada e relatórios de baixa acionabilidade e indicam as dimensões em que abordagens baseadas em IA poderiam agregar valor. É justamente esse potencial de complementaridade que a seção seguinte examina, do lado oposto do espectro: o que os LLMs, isoladamente, conseguem detectar quando confrontados com a mesma tarefa.


\section{LLMs para detecção de vulnerabilidades em código}\label{sec:llms-deteccao-vulnerabilidades}
Se a Seção Seção~\ref{sec:avaliacoes-sast} caracterizou o braço SAST do experimento, esta seção caracteriza o material bruto do segundo braço, sendo a capacidade de LLMs de propósito geral em identificar vulnerabilidades , que mais adiante será encapsulado, no caso deste TCC, pelas ferramentas empresariais de desenvolvimento assistido por IA (Seção ~\ref{sec:agentes-ia-seguranca}). Entender onde os LLMs falham quando avaliados diretamente é pré-condição para interpretar corretamente o desempenho dessas ferramentas quando operam como agentes sobre um repositório. A literatura sobre LLMs para detecção de vulnerabilidades cresceu rapidamente: revisões sistemáticas recentes catalogam de 58 a 208 estudos primários \cite{Sheng2025LLMSurvey, Zhou2025LLMReview, Germano2025LLMSLR}. O quadro geral é de potencial acompanhado de fragilidades. No maior estudo empírico com modelos de propósito geral, 16 LLMs avaliados sobre 5 mil trechos de código atingiram acurácia média de 62,8\% e F1 médio de 0,71, com queda de cerca de 10 pontos percentuais ao passar de dados sintéticos para dados reais \cite{Khare2025LLMEffectiveness}. Avaliações mais rigorosas são ainda menos otimistas: quando benchmarks tradicionais foram corrigidos quanto a rótulos errados e vazamento de dados, o F1 de modelos de código despencou de 68,3\% para 3,1\%, e nem o GPT-4 com raciocínio em cadeia superou o acaso na distinção entre função vulnerável e sua versão corrigida \cite{Ding2025CodeLMVuln}. Estudos de raciocínio confirmam o diagnóstico: LLMs exibem consistência semântica inferior a 50\% em tarefas estruturadas de análise de vulnerabilidades \cite{Li2025SVTrustEval}, erram a classificação da CWE específica mesmo quando "percebem" o problema (cerca de 16\% de acerto, próximo do aleatório) \cite{Gaspar2025OpenLLM, Lin2025LLMCompare} e degradam a localização em arquivos longos \cite{Sovrano2025LostInEnd}.

Um segundo eixo de resultados diz respeito ao papel do contexto. Avaliações com trechos isolados de código subestimam ou distorcem a capacidade dos modelos; com contexto adequado do repositório, a acurácia estrita chegou a 67\% no benchmark CORRECT, e a maioria dos falsos positivos foi atribuída a erros de raciocínio, não de classificação \cite{Li2025LLMVDAsk}. Na escala de projeto, porém, detectores baseados em LLM apresentaram \textit{recall} entre 21\% e 34\% com taxas de falsos alertas de 85–97\%, custo computacional elevado e causas de erro dominadas por análise interprocedural rasa \cite{Li2026ProjectScaleLLM}. Abordagens que fornecem contexto estruturado ao modelo  fatias de código guiadas por grafos de propriedade \cite{Lekssays2025LLMxCPG}, abstrações de funções chamadas \cite{Yang2025ContextVD} ou especificações extraídas de correções históricas \cite{Zhu2025SpecGuidedVD} , melhoram consistentemente os resultados, o que sugere que a forma de entregar o código ao modelo importa tanto quanto o modelo em si. Essa observação é diretamente relevante para o TCC: Cursor e Codex operam com camadas de contexto e ferramentas sobre o repositório \cite{Cursor2025Docs, OpenAI2025CodexDocs}; isso os distingue dos LLMs isolados revisados nesta seção, sem permitir presumir a priori superioridade de desempenho.

Uma terceira vertente investiga modelos especializados por \textit{fine-tuning}. Para Python especificamente, o SecureQwen reportou F1 entre 84\% e 99\% por categoria CWE, embora com rótulos gerados por um \textit{ensemble} de ferramentas SAST (incluindo Bandit e Semgrep), o que introduz circularidade na avaliação \cite{Mechri2025SecureQwen}. Resultados similares de alto desempenho em benchmarks controlados foram reportados para C/C++ \cite{Ferrag2025SecureFalcon}, mas os ganhos do \textit{fine-tuning} sobre linhas de base são frequentemente modestos em dados desbalanceados \cite{Shestov2025LLMFinetune}. Cabe notar que esses modelos dedicados diferem substancialmente das ferramentas empresariais de desenvolvimento assistido por IA avaliadas neste trabalho, que utilizam modelos generalistas com acesso a contexto de repositório. Por fim, a dimensão da explicação permanece pouco explorada: apenas 5,3\% dos 208 estudos de uma revisão sistemática abordam explicação de vulnerabilidades, contra 91,3\% dedicados à detecção binária \cite{Germano2025LLMSLR}; o trabalho que propôs rubricas de corretude, clareza e acionabilidade constatou que alta taxa de detecção não implica explicações acionáveis \cite{Mao2025ExplainableVD}. Tomados em conjunto, os dois primeiros eixos revisados nesta seção , desempenho isolado modesto e sensibilidade forte ao contexto motivam diretamente a pergunta central da próxima seção: quando SAST e LLM são colocados lado a lado sob o mesmo protocolo, qual abordagem efetivamente prevalece, e sob quais condições?


\section{Comparações entre SAST tradicional e abordagens baseadas em IA}\label{sec:comparacoes-sast-ia}
As comparações diretas entre SAST e LLMs formam a linha de pesquisa mais próxima deste TCC: é aqui que o \textit{trade-off} entre os dois braços do experimento é medido explicitamente, sob o mesmo conjunto de dados. O estudo pioneiro comparou o ChatGPT com onze ferramentas SAST para PHP e encontrou o padrão que se repetiria nos trabalhos seguintes: o LLM detectou 62–68\% das vulnerabilidades, contra 32\% da melhor ferramenta tradicional, mas alertou incorretamente 91\% dos trechos seguros e variou 6 pontos percentuais entre execuções idênticas \cite{Ozturk2023AIvsSAST}. Avaliação posterior com cinco LLMs no mesmo domínio confirmou taxas de detecção de até 61,5\% acompanhadas de falsos positivos superiores a 60\% em código corrigido \cite{Cetin2024LLMPHP}. Estudos anteriores com modelos da geração GPT-3 já apontavam a mesma assimetria ,\textit{recall} elevado com taxa de falsos positivos várias vezes superior à das ferramentas estáticas \cite{Purba2023LLMVuln}.

Comparações metodologicamente mais rigorosas refinaram o quadro. Em benchmark controlado sobre 578 arquivos Java, o melhor LLM naquele momento (GPT-4 com raciocínio em cadeia, prompt CoT de 8 passos) atingiu F1 de 0,672 e cobriu 13 das 17 categorias CWE, contra F1 de 0,655 e cobertura de 11 categorias da melhor ferramenta SAST; em contrapartida, o SAST produziu menos falsos positivos (23 contra 55), com custo financeiro nulo e execução em menos de um minuto \cite{Tamberg2025LLMHarness}. Em dez projetos C\# reais, LLMs de fronteira superaram três ferramentas SAST em F1 (0,75–0,80 contra 0,26–0,55), mas falharam completamente na localização precisa (linha/coluna) dos achados \cite{Gnieciak2025LLMvsSAST}. Um estudo de larga escala com 15 ferramentas SAST e 12 LLMs de código aberto sobre repositórios reais em Java, C e Python — um dos poucos a incluir Python nesta escala , reportou o mesmo padrão qualitativo já visto acima (SAST com baixa cobertura e poucos falsos positivos; LLMs com cobertura muito superior e falsos positivos elevados, mitigados apenas parcialmente por \textit{ensemble}), embora se trate de preprint ainda sem revisão por pares quase dois anos após a submissão, o que recomenda cautela redobrada quanto aos valores absolutos reportados \cite{Zhou2024SASTLLMRepo}.

No único benchmark identificado com foco em Python e revisão por pares pendente, cobrindo 26 repositórios e 796 entradas rotuladas, os \textit{scanners} baseados em LLMs de propósito geral superaram amplamente as ferramentas de regras , o Semgrep obteve F3 de 17,7, contra 51,7 de um LLM comercial e 73,0 da melhor ferramenta especializada \cite{Pellew2026RealVuln}. Também em Python, um estudo comparando aprendizado de máquina, ChatGPT e análise estática sobre 21 scripts concluiu que as abordagens baseadas em IA superaram significativamente as ferramentas estáticas, embora a amostra reduzida limite a generalização \cite{Farasat2025PythonSecurity}. Um preprint mais recente (janeiro de 2026) propõe um protocolo de avaliação hierárquica por CWE especificamente para Bandit, Semgrep e CodeQL contra LLMs abertos em nível de função em Python, reforçando que Python começa a receber atenção comparativa dedicada, ainda que os resultados numéricos desse trabalho não tenham podido ser verificados de forma independente nesta revisão e, por isso, não sejam aqui reproduzidos \cite{Adnan2026ALPHA}.


\begin{table}[htbp]
\centering
\caption{Estudos comparativos SAST–LLM mais próximos do presente trabalho.}\label{tab:compartivos-sast-llm}
\begin{tabularx}{\textwidth}{|X|X|X|X|X|}
\hline
\textbf{Estudo} & \textbf{Linguagem e dados} & \textbf{Comparação} & \textbf{Métricas e achados selecionados} & \textbf{Situação} \\
\hline
Ozturk et al. \cite{Ozturk2023AIvsSAST} & PHP & ChatGPT vs. 11 SAST & Detecção de 62–68\% vs. 32\%; 91\% de falsos alertas em trechos seguros; variação de 6 p.p. & Estudo pioneiro \\
\hline
Tamberg et al. \cite{Tamberg2025LLMHarness} & Java, 578 arquivos & GPT-4 com raciocínio em cadeia vs. SAST & F1 de 0,672 vs. 0,655; 13 vs. 11 CWEs; 55 vs. 23 falsos positivos & Revisado por pares \\
\hline
Gnieciak e Szandala \cite{Gnieciak2025LLMvsSAST} & C\#, 10 projetos & LLMs de fronteira vs. 3 SAST & F1 de 0,75–0,80 vs. 0,26–0,55; falha de localização precisa pelos LLMs & Preprint \\
\hline
Zhou et al. \cite{Zhou2024SASTLLMRepo} & Java, C e Python; repositórios & 15 SAST vs. 12 LLMs & Maior cobertura dos LLMs, com maior proporção de funções marcadas & Preprint \\
\hline
Pellew e Raza \cite{Pellew2026RealVuln} & Python, 26 repositórios, 796 entradas & Ferramentas de regras vs. \textit{scanners} LLM & F3: Semgrep 17,7; LLM comercial 51,7; melhor ferramenta especializada 73,0 & Preprint \\
\hline
Farasat e Posegga \cite{Farasat2025PythonSecurity} & Python, 21 scripts & Aprendizado de máquina, ChatGPT e SAST & Abordagens de IA superaram SAST; amostra reduzida & Revisado por pares \\
\hline
\end{tabularx}f
\end{table}

A Tabela ~\ref{tab:compartivos-sast-llm} torna explícita a heterogeneidade da evidência: os trabalhos diferem em linguagem, granularidade, conjunto de dados e métrica. Por essa razão, não é metodologicamente apropriado agregar seus números em um gráfico único; a tabela preserva a comparabilidade limitada e evita inferências que os estudos não sustentam.

Em síntese, a literatura comparativa sugere um \textit{trade-off} consistente: LLMs oferecem maior cobertura e explicações em linguagem natural, ao custo de falsos positivos elevados, não-determinismo, custo financeiro e localização imprecisa; ferramentas SAST oferecem precisão, reprodutibilidade e custo baixo, ao custo de cobertura limitada \cite{Ozturk2023AIvsSAST, Tamberg2025LLMHarness, Gnieciak2025LLMvsSAST}. Nenhum desses estudos, contudo, avaliou ferramentas empresariais de desenvolvimento assistido por IA operando como agentes sobre o repositório ,cenário deste TCC , e mesmo somando os trabalhos mais recentes localizados nesta revisão, o ecossistema Python permanece coberto majoritariamente por preprints \cite{Pellew2026RealVuln, Farasat2025PythonSecurity, Zhou2024SASTLLMRepo, Adnan2026ALPHA}. Diante de um \textit{trade-off} tão bem documentado entre as duas abordagens, é natural que uma parcela crescente da literatura tenha deixado de tratá-las como alternativas excludentes e passado a investigar formas de combiná-las , é o que a seção seguinte revisa.




\section{Abordagens híbridas SAST + IA}\label{sec:abordagens-hibridas}
Diante do \textit{trade-off} descrito na seção anterior, uma linha crescente de trabalhos combina análise estática e LLMs , o terceiro braço deste TCC, e o que a literatura sugere ser o de maior potencial prático, ainda que o mais carente de validação padronizada. Três padrões arquiteturais predominam. No primeiro o mais próximo do desenho deste TCC , o LLM atua a jusante do SAST, validando, filtrando e explicando os alertas. Um sistema multi-agente que revisa candidatos do Semgrep com contexto do código reduziu a taxa de falsos positivos de 0,423 para 0,048 (redução de 88,6\%) ao custo de apenas 3,1\% dos verdadeiros positivos, elevando o F1 de 0,784 para 0,909 no OWASP Benchmark \cite{Ameen2026QASecClaw}. Resultados da mesma ordem foram reportados ao tratar saídas SAST como contratos estruturados antes da adjudicação por LLM (F1 de até 0,955 em projetos reais) \cite{Iranmanesh2025ZeroFalse}, na triagem de relatórios com raciocínio LLM (precisão de 9,4\% para 52,1\%) \cite{Li2025SASTReportsLLM} e na integração Semgrep + LLM com redução de aproximadamente 91\% dos alertas \cite{Agrawal2025SASTGenius}. Na comparação experimental de cenários, o híbrido SAST + LLM superou ambas as abordagens isoladas nos cenários de menor criticidade, embora o SAST puro tenha permanecido preferível quando o \textit{recall} é prioritário, e o híbrido tenha inclusive piorado o desempenho em categorias específicas \cite{Ciavatta2026LLMSAST}.

No segundo padrão, o LLM atua a montante, ampliando a capacidade do motor estático: inferindo especificações de \textit{taint} (\textit{sources} e \textit{sinks}) para o CodeQL, o que dobrou a detecção de CVEs reais (de 22,5\% para 45,8\%) \cite{Li2025IRIS}; descobrindo \textit{sources/sinks/sanitizers} em escala industrial \cite{Luo2025SASTLLMDFA}; gerando consultas CodeQL a partir de descrições de CWE \cite{Wang2026CQLLM}; ou sintetizando regras de forma autônoma com validação por prova de conceito, abordagem que descobriu CVEs novos em pacotes Python do PyPI \cite{Tsigkourakos2026QRS}. O terceiro padrão emprega o SAST como oráculo de realimentação em laços de geração e correção: em Python, o retorno de alertas do Bandit ao LLM reduziu o código vulnerável gerado de cerca de 40\% para menos de 10\% \cite{Alrashedy2025FDSP, Lee2026CodeEnhancer}. Um estudo independente, em escala de produção (733 trechos de código Python e JavaScript gerados por GitHub Copilot, CodeWhisperer e Codeium extraídos de projetos reais do GitHub, com 29,5\% dos trechos Python contendo fraquezas de segurança em 43 categorias de CWE), confirma esse mesmo padrão de realimentação SAST→LLM fora de ambiente controlado: fornecer ao Copilot Chat os alertas gerados por ferramentas de análise estática permitiu corrigir até 55,5\% das fraquezas identificadas — evidência revisada por pares, em Python, de que uma ferramenta comercial de IA pode atuar como agente corretor quando municiada por sinal de SAST, mesmo que o artigo não avalie a mesma ferramenta como agente \textit{detector} \cite{Fu2025CopilotSecurity}.

A leitura crítica desses resultados exige três ressalvas. Primeiro, boa parte dos trabalhos híbridos mais próximos do TCC são preprints ainda não revisados por pares \cite{Ameen2026QASecClaw, Iranmanesh2025ZeroFalse, Agrawal2025SASTGenius, Tsigkourakos2026QRS}. Segundo, os experimentos concentram-se em Java (frequentemente no OWASP Benchmark sintético) e Solidity, com escassa validação em Python \cite{Ciavatta2026LLMSAST, Ameen2026QASecClaw, Li2025SASTReportsLLM} — o estudo de Fu et al. é uma das poucas exceções revisadas por pares com dados reais de Python, ainda que restrito ao padrão de correção, não de detecção \cite{Fu2025CopilotSecurity}. Terceiro, a combinação não é garantidamente benéfica: a união ingênua de detectores baseados em regras e em aprendizado piora o desempenho global \cite{Croft2021RuleLearningSAST}, e a filtragem por LLM pode degradar categorias específicas de vulnerabilidade \cite{Ciavatta2026LLMSAST, Ameen2026QASecClaw}. Surveys dedicados à integração LLM + análise estática convergem, ainda assim, para a conclusão de que a combinação supera cada abordagem isolada na maioria dos cenários, posicionando o paradigma híbrido como direção promissora, porém carente de validação empírica padronizada \cite{Jin2025LLMSASTSurvey, Khare2025LLMEffectiveness}. Falta a essa linha de pesquisa, contudo, um ingrediente específico: os padrões híbridos revisados aqui usam LLMs via API, orquestrados por \textit{pipelines} construídos pelos próprios autores, nenhum emprega diretamente uma ferramenta empresarial de desenvolvimento assistido por IA operando de forma autônoma e agêntica sobre um repositório, o que remete à seção seguinte.


\section{Agentes de IA e workflows automatizados para segurança de código}\label{sec:agentes-ia-seguranca}
Esta seção trata do paradigma do qual emergem as ferramentas avaliadas neste TCC, sendo eles Cursor e OpenAI Codex: agentes de IA comerciais que operam em ciclos de raciocínio e ação sobre um repositório, com acesso a ferramentas. Claude Code é mencionado somente como alternativa comercial contextual, não como objeto de avaliação. Entender o que a literatura acadêmica já demonstrou (e o que ainda não demonstrou) sobre esse paradigma aplicado à segurança de código é o elo final antes de justificar formalmente o desenho experimental na Seção ~\ref{sec:lacunas-literatura}. O paradigma foi consolidado em engenharia de software por benchmarks de resolução de \textit{issues} reais \cite{Jimenez2024SWEbench} e pela demonstração de que interfaces agente-computador especializadas multiplicam o desempenho dos modelos \cite{Yang2024SWEagent}. Cursor e Codex implementam esse paradigma com modos de agente, acesso ao repositório e ferramentas de execução configuráveis \cite{Cursor2025Docs, OpenAI2025CodexDocs}. Sua transposição para segurança de código, porém, é recente e apresenta resultados mistos. No benchmark JITVUL, agentes do tipo ReAct com acesso a ferramentas melhoraram a acurácia em pares em relação ao \textit{prompt} simples em alguns cenários, mas o desempenho absoluto permaneceu baixo (inferior a 21\%) e o simples enriquecimento do \textit{prompt} com contexto de dependências por vezes superou o agente completo \cite{Yildiz2025LLMAgents}.

Arquiteturas multi-agente mais elaboradas mostram resultados promissores em cenários específicos. Um \textit{ensemble} de agentes orquestrando análise estática com recuperação de conhecimento de vulnerabilidades encontrou falhas em bases de código industriais nas quais o CodeQL isolado nada detectou, com custo da ordem de poucos dólares por repositório \cite{Liang2026Argus}; a tríade neuro-simbólica de agentes que sintetiza e valida regras descobriu cinco CVEs novos em pacotes populares do PyPI \cite{Tsigkourakos2026QRS}; e agentes de filtragem sustentam os melhores resultados de redução de falsos positivos \cite{Ameen2026QASecClaw}. Um exemplo particularmente informativo, publicado ems veículo revisado por pares de primeira linha (ICSE 2026), é o VulTrial: uma arquitetura multi-agente inspirada em um julgamento (agentes com papéis de pesquisador de segurança, autor do código, moderador e conselho de revisão) que, sobre o benchmark PrimeVul-pair (C/C++), quase dobrou o número de pares corretamente classificados em relação a \textit{baselines} single-agent e multi-agente simples (81 contra 40–44 pares corretos com GPT-4o, chegando a 96 com ajuste do agente moderador) e identificou três vulnerabilidades zero-day confirmadas em projetos reais em produção e reportando,  o custo em dólares de cada configuração (de US\$ 7,46 com GPT-3.5 a cerca de US\$ 23,90 com GPT-4o para o conjunto de teste completo), uma das poucas evidências quantitativas de custo disponíveis nesta linha de pesquisa \cite{Widyasari2026VulTrial}. Ainda assim, o VulTrial usa LLMs via API orquestrados por uma arquitetura de agentes construída pelos próprios autores , não Cursor, Codex ou Claude Code e é validado exclusivamente em C/C++, não em Python.

No extremo oposto da cadeia, o único estudo de campo com ferramenta de detecção e reparo integrada ao IDE encontrou utilidade percebida baixa (2,5/5 para alertas) e 75\% de correções não aplicáveis, evidenciando a distância entre benchmark e prática \cite{Steenhoek2025IDEUserStudy}. Uma busca complementar dedicada especificamente a essa lacuna localizou um preprint que avalia o recurso "Code Review" do GitHub Copilot; seu achado qualitativo é que a ferramenta falha em identificar vulnerabilidades críticas e concentra comentários em problemas de estilo \cite{Amro2025CopilotReview}. Há também evidência acadêmica direta, porém adjacente, para Cursor: o SecCodePLT avaliou manualmente 153 exemplos Python de geração e completamento seguro e observou taxas de código seguro de 52,8\% em testes dinâmicos e 62,0\% em testes baseados em regras, sem política de segurança; com política explícita, os valores subiram para 67,4\% e 86,7\%, respectivamente \cite{Nie2025SecCodePLT}. Esse estudo avalia a segurança do código que Cursor gera, e não sua capacidade de detectar vulnerabilidades já existentes em um repositório. A Tabela ~\ref{tab:ferramentas-empresariais-estudo} delimita essas evidências próximas e evita confundir geração segura, revisão de \textit{pull request} e detecção de vulnerabilidades.



\begin{table}[htbp]
\centering
\caption{Ferramentas empresariais de IA em estudos de segurança adjacentes ao TCC}\label{tab:ferramentas-empresariais-estudo}
\begin{tabularx}{\textwidth}{|X|X|X|X|X|}
\hline
\textbf{Ferramenta ou sistema} & \textbf{Estudo} & \textbf{Cenário} & \textbf{Métrica ou achado} & \textbf{Relação com o TCC} \\
\hline
Cursor & Nie et al. \cite{Nie2025SecCodePLT} & 153 exemplos Python; geração e completamento seguro & 52,8\% em testes dinâmicos e 62,0\% em testes de regras; políticas de segurança elevaram os valores para 67,4\% e 86,7\% & Avaliação direta de Cursor em segurança, mas não de detecção \\
\hline
GitHub Copilot Code Review & Amro e Alalfi \cite{Amro2025CopilotReview} & Revisão de código vulnerável & Falhas qualitativas em vulnerabilidades críticas; sem métricas verificadas nesta revisão & Produto comercial adjacente; preprint \\
\hline
Protótipo integrado ao IDE & Steenhoek et al. \cite{Steenhoek2025IDEUserStudy} & Estudo de campo de detecção e reparo & Utilidade dos alertas: 2,5/5; 75\% das correções não aplicáveis & Não avalia Cursor ou Codex \\
\hline
Claude Code & Luo et al. \cite{Luo2026CodeAugur} & \textit{Baseline} em benchmark de agentes com validação dinâmica & Comparado a um agente acadêmico de inferência de especificações e \textit{fuzzing} & Alternativa comercial contextual; preprint e fora de Python \\
\hline
\end{tabularx}
\end{table}

Portanto, as evidências encontradas não avaliam Cursor ou OpenAI Codex, configurados como ferramentas-agente de desenvolvimento, para detectar vulnerabilidades em aplicações Python. Elas reforçam a relevância do tema, mas tratam de geração segura, revisão de código, protótipos de IDE ou agentes acadêmicos em outros contextos.


\section{Lacunas identificadas na literatura}\label{sec:lacunas-literatura}
A análise consolidada nas seções anteriores permite identificar as seguintes lacunas, que fundamentam o experimento deste TCC. Esta síntese incorpora os resultados de uma busca complementar dedicada, realizada em julho de 2026 com mais de dez consultas distintas em bases revisadas por pares (ACM DL, IEEE Xplore, Springer, USENIX, NDSS) e complementadas por buscas gerais na web, especificamente voltada a verificar se alguma das lacunas abaixo já havia sido preenchida por literatura publicada entre o fechamento da revisão original e a data desta revisão.

\textbf{Ausência de comparação direta entre SAST tradicional, ferramentas-agente empresariais e híbrido em um mesmo protocolo.} Os estudos comparativos opõem SAST a LLMs \cite{Ozturk2023AIvsSAST, Tamberg2025LLMHarness, Gnieciak2025LLMvsSAST, Zhou2024SASTLLMRepo} ou avaliam híbridos contra suas linhas de base \cite{Ciavatta2026LLMSAST, Ameen2026QASecClaw}, mas nenhum trabalho identificado , nem mesmo entre os localizados na busca complementar de 2026 , compara as três abordagens simultaneamente, sob o mesmo conjunto de dados e as mesmas métricas.

\textbf{Sub-representação do ecossistema Python.} Um mapeamento com 180 estudos de predição de vulnerabilidades encontrou apenas 2 dedicados a Python, contra 86 para C/C++ \cite{Kalouptsoglou2023VulnPrediction}; em revisão com 208 estudos sobre LLMs, Python representa 7,7\% \cite{Germano2025LLMSLR}. A busca complementar de 2026 confirma que a lacuna está diminuindo, mas não fechando: os trabalhos localizados com foco em Python continuam sendo, quase sem exceção, preprints \cite{Pellew2026RealVuln, Marquer2026SAGA, Tsigkourakos2026QRS, Zhou2024SASTLLMRepo, Adnan2026ALPHA} ou usam amostras muito reduzidas \cite{Farasat2025PythonSecurity}; a única exceção revisada por pares com dados reais de Python em escala (Fu et al., ACM TOSEM) trata de geração e correção de código, não de detecção de vulnerabilidades pré-existentes \cite{Fu2025CopilotSecurity}.

\textbf{Ferramentas empresariais de desenvolvimento assistido por IA não avaliadas como detectores.} Esta é a lacuna mais crítica identificada neste levantamento. Não foi localizado estudo revisado por pares que avalie Cursor ou OpenAI Codex, configurados como ferramentas-agente de desenvolvimento, para detecção de vulnerabilidades em aplicações Python. As evidências mais próximas são agentes de pesquisa construídos com LLMs via API \cite{Yildiz2025LLMAgents, Widyasari2026VulTrial}, um protótipo de detecção e reparo integrado ao IDE \cite{Steenhoek2025IDEUserStudy}, revisão de código por GitHub Copilot \cite{Amro2025CopilotReview} e avaliação de Cursor para geração segura de código \cite{Nie2025SecCodePLT}. Esses estudos reforçam a relevância da comparação, mas não substituem a avaliação dos dois produtos-alvo como detectores de vulnerabilidades existentes em repositórios Python.

\textbf{Reprodutibilidade, não-determinismo e dependência de prompt.} LLMs variam entre execuções idênticas \cite{Ozturk2023AIvsSAST}, seus resultados são fortemente sensíveis à estratégia de \textit{prompt} com ganhos de até 0,18 em F1 entre estratégias \cite{Khare2025LLMEffectiveness} e diferenças superiores a 0,3 em F1 entre formulações de instrução \cite{Iranmanesh2025ZeroFalse} , e estratégias supostamente superiores, como \textit{few-shot}, podem degradar o desempenho \cite{Lin2025LLMCompare}. Protocolos padronizados de avaliação, como \textit{harnesses} dedicados, são resposta ainda incipiente \cite{Tamberg2025LLMHarness}.

\textbf{Fragilidade dos benchmarks.} Vazamento de dados, rótulos incorretos e distribuições irreais inflam os resultados reportados \cite{Ding2025CodeLMVuln, Chen2023DiverseVul, Chakraborty2022DLLimits}; avaliações em nível de função superestimam a capacidade em relação ao nível de repositório \cite{Wen2026VulEval}. A própria dispersão de resultados observada na busca complementar de 2026 , por exemplo, a faixa de 90–100\% de detecção por LLM reportada em \cite{Zhou2024SASTLLMRepo}, atipicamente alta frente aos 62,8\% médios de \cite{Khare2025LLMEffectiveness} ilustra esse risco e reforça a necessidade, no experimento do TCC, de conjuntos de dados verificados e métricas resistentes (avaliação em código seguro, penalização de falsos positivos).

\textbf{Custo e limites operacionais das ferramentas de IA.} As avaliações que reportam custo indicam despesas e tempos ordens de magnitude superiores aos do SAST \cite{Tamberg2025LLMHarness, Li2026ProjectScaleLLM}, e o VulTrial acrescenta um dado ponto de referência recente (US\$ 7,46 a US\$ 23,90 por execução completa de um conjunto de teste de 435 pares, a depender do modelo e da arquitetura de agentes) \cite{Widyasari2026VulTrial}. Ainda assim, essa dimensão permanece rara nas comparações, e a busca complementar de 2026 não localizou nenhum estudo acadêmico revisado por pares que reporte custo ou limites de uso (\textit{rate limits}, cotas) especificamente para Cursor, OpenAI Codex ou Claude Code aplicados a detecção de vulnerabilidades apenas material de mercado (documentação de precificação, relatos de usuários), inadequado como evidência acadêmica e por isso não incorporado a este capítulo.

\textbf{Pouca avaliação da utilidade prática dos relatórios.} Apenas uma pequena parte dos estudos abordam explicação \cite{Germano2025LLMSLR}; a indústria demanda diagnósticos acionáveis \cite{Li2025IndustrySAST}; e a utilidade percebida em campo é baixa mesmo quando as métricas de benchmark são razoáveis \cite{Steenhoek2025IDEUserStudy}. Rubricas de corretude, clareza e acionabilidade existem \cite{Mao2025ExplainableVD}, mas raramente são aplicadas em comparações entre abordagens.

Em conjunto, essas lacunas justificam o desenho deste trabalho: uma comparação empírica, sob protocolo único, entre SAST tradicional (Bandit e Semgrep Community Edition), Cursor e OpenAI Codex como ferramentas-agente empresariais, e uma abordagem híbrida SAST + IA, aplicada a código Python, combinando métricas quantitativas (precisão, \textit{recall}, F1, taxa de falsos positivos, tempo e custo) e avaliação qualitativa da clareza, da explicação e da priorização dos relatórios produzidos. A busca complementar realizada especificamente para esta revisão reforça a urgência desse desenho: apesar da expansão de agentes de codificação comerciais e de estudos adjacentes de segurança, não foi localizado estudo revisado por pares que avalie Cursor ou OpenAI Codex como detectores de vulnerabilidades em aplicações Python , o espaço específico que este TCC ocupa.

\chapter{Metodologia}\label{cap:metodologia}

Este capítulo operacionaliza o problema de pesquisa delimitado nos capítulos anteriores: comparar a capacidade de ferramentas SAST tradicionais, de agentes empresariais de codificação baseados em IA e de uma abordagem híbrida SAST + IA para detectar, explicar e priorizar vulnerabilidades em aplicações Python. A fundamentação e os critérios de escolha de Bandit, Semgrep Community Edition, Cursor e OpenAI Codex foram apresentados no Capítulo~\ref{cap:fundTeoric}; neste capítulo, o foco recai sobre como essas ferramentas serão executadas sobre o benchmark RealVuln, quais dados serão coletados e como os resultados serão avaliados.

A metodologia adota um desenho empírico comparativo, com medidas repetidas sobre os mesmos repositórios Python e o \textit{ground truth} público e versionado do RealVuln \cite{Pellew2026RealVuln}. Um repositório controlador organizará a versão fixada do benchmark, as cópias entregues às ferramentas, as instruções e os resultados, mantendo os arquivos de resposta fora do alcance dos agentes durante a análise. Bandit e Semgrep serão invocados pelo \textit{harness} em contêineres Docker efêmeros e equivalentes. As ferramentas empresariais de IA serão avaliadas sobre cópias idênticas dos mesmos repositórios, com instruções, permissões e recursos controlados.

\section{Questões de Pesquisa}\label{sec:questoes-pesquisa-metodologia}

O experimento é orientado pelas seguintes questões de pesquisa:

\begin{enumerate}
\item[\textbf{QP1}] Qual é o desempenho de cada condição experimental na detecção de vulnerabilidades Python, considerando precisão, \textit{recall}, F1, F2, F3 e taxa de falsos positivos?
\item[\textbf{QP2}] Qual é a qualidade dos relatórios verdadeiros-positivos quanto à correção técnica, localização, clareza e acionabilidade?
\item[\textbf{QP3}] Em que medida fornecer os alertas de Bandit e Semgrep aos agentes de IA melhora ou degrada a detecção e a qualidade dos relatórios em relação às ferramentas executadas isoladamente?
\item[\textbf{QP4}] Quão estáveis são os resultados das condições baseadas em IA entre execuções independentes e quão adequadamente essas condições indicam a prioridade dos achados quando existe uma referência de severidade confiável?
\end{enumerate}

As QP1 e QP3 constituem o eixo quantitativo principal. A QP2 acrescenta a avaliação da utilidade prática dos relatórios, enquanto a QP4 combina medidas quantitativas de estabilidade e concordância com a análise da priorização, dimensões apontadas como pouco exploradas pela literatura revisada \cite{Germano2025LLMSLR, Mao2025ExplainableVD, Tamberg2025LLMHarness}. O tempo de execução e o consumo de tokens serão apresentados como caracterização operacional complementar, sem constituir uma questão de pesquisa autônoma.

\section{Classificação e Desenho da Pesquisa}\label{sec:classificacao-pesquisa}

Quanto à natureza, a pesquisa é aplicada, pois produz evidência empírica destinada a apoiar a escolha e a combinação de ferramentas de análise de segurança no ecossistema Python. Quanto à abordagem, emprega métodos mistos com predominância quantitativa: as condições são comparadas por métricas derivadas da matriz de confusão, enquanto a qualidade dos relatórios é examinada por uma rubrica estruturada. Quanto aos objetivos, o estudo é exploratório e descritivo-comparativo. Não se pretende demonstrar causalidade geral sobre todas as ferramentas ou modelos de IA, mas caracterizar diferenças observadas sob as condições controladas deste protocolo.

O procedimento central é um experimento computacional de medidas repetidas: cada condição experimental é aplicada aos mesmos repositórios da versão fixada do RealVuln. A comparação utiliza as vulnerabilidades e os falsos-positivos armadilha já rotulados pelo benchmark, reduzindo a influência de diferenças de corpus entre as ferramentas. A pesquisa bibliográfica apresentada nos capítulos anteriores fundamenta a escolha das condições, dos dados, das métricas e das ameaças à validade; o desenvolvimento do \textit{harness} constitui a instrumentação necessária para integrar Cursor e Codex ao formato de resultados e ao avaliador do benchmark.

\section{Condições Experimentais}\label{sec:condicoes-experimentais}

As três abordagens conceituais do trabalho são operacionalizadas em seis condições, apresentadas na Tabela~\ref{tab:condicoes-experimentais}. Bandit e Semgrep são avaliados separadamente, conforme suas diferenças técnicas discutidas na Seção~\ref{subsec:ferramentas-sast}. Cursor e Codex também permanecem separados, evitando que uma união artificial de suas respostas seja interpretada como uma única ferramenta. No braço híbrido, a união normalizada e sem duplicatas dos alertas de Bandit e Semgrep é fornecida, como contexto adicional, a uma nova execução de cada agente.

\begin{table}[htb]
\ABNTEXfontereduzida
\caption{\label{tab:condicoes-experimentais}Condições experimentais avaliadas}
\centering
\begin{tabularx}{\textwidth}{|l|X|X|}
\hline
\textbf{Condição} & \textbf{Execução} & \textbf{Papel na comparação} \\
\hline
C1 & Bandit em configuração fixada & Linha de base SAST específica para Python \\
\hline
C2 & Semgrep Community Edition com conjunto de regras fixado & Linha de base SAST baseada em regras declarativas \\
\hline
C3 & Cursor com GPT-5.6 Luna, sem receber alertas SAST & Agente empresarial isolado com modelo controlado \\
\hline
C4 & OpenAI Codex com GPT-5.6 Luna, sem receber alertas SAST & Agente empresarial isolado com modelo controlado \\
\hline
C5 & Alertas normalizados de C1 e C2 fornecidos ao Cursor com GPT-5.6 Luna & Condição híbrida SAST $\rightarrow$ Cursor com modelo controlado \\
\hline
C6 & Alertas normalizados de C1 e C2 fornecidos ao Codex com GPT-5.6 Luna & Condição híbrida SAST $\rightarrow$ Codex com modelo controlado \\
\hline
\end{tabularx}
\legend{Fonte: Elaborado pelo autor (2026).}
\end{table}

Os resultados primários serão reportados por condição, e não apenas agregados nos rótulos ``SAST'', ``IA'' e ``híbrido''. Uma síntese por abordagem poderá ser apresentada como leitura complementar, desde que preserve os valores individuais e não trate ferramentas diferentes como observações independentes de uma população de ferramentas.

O bloco experimental principal utilizará o GPT-5.6 Luna em C3--C6, desde que o modelo permaneça disponível nos dois produtos no início da coleta. A adoção do mesmo modelo reduz o confundimento entre a capacidade do modelo subjacente e os efeitos próprios da ferramenta, como seleção de contexto, orquestração, ferramentas disponíveis e interface agêntica. O identificador do modelo, o nível de raciocínio, o modo de contexto e as demais opções observáveis serão fixados no estudo-piloto e mantidos durante todo o bloco \cite{Cursor2026Models, OpenAI2026ModelGuidance}.

Concluído o bloco principal, o desenho admite um bloco complementar com as configurações nativas indicadas pelos fornecedores. Na data de planejamento, os candidatos são o modo Auto no Cursor e o GPT-5.6 Terra no Codex. O Auto é uma estratégia adaptativa de seleção e pode encaminhar solicitações a modelos diferentes; por isso, não será tratado como equivalente ao Terra nem combinado com a linha de base Luna. A disponibilidade desse bloco será decidida no estudo-piloto, antes da observação dos resultados principais, a partir das cotas e do limite de tokens reservados para a coleta. Se executado, abrangerá o corpus completo e seguirá as mesmas três repetições, instruções, permissões e critérios de avaliação. Caso os recursos disponíveis não comportem o bloco completo, ele não será executado, evitando uma seleção posterior de repositórios ou resultados favoráveis. A configuração efetiva, a justificativa documental e a data de consulta serão registradas antes de iniciar o bloco.

Nas condições C3 e C4, os agentes poderão ler e pesquisar os arquivos do repositório, preservando a característica agêntica definida na Seção~\ref{sec:agentes-codificacao}, mas não poderão executar Bandit, Semgrep ou outros analisadores de segurança. Nas condições C5 e C6, receberão os alertas SAST previamente coletados, sem acesso ao \textit{ground truth}. O acesso à rede, a execução de comandos, a possibilidade de modificar arquivos e as demais permissões serão configurados da forma mais equivalente permitida pelos dois produtos e registrados como parte do artefato experimental.

\section{Universo, Corpus e Unidade de Análise}\label{sec:universo-amostra}

O universo de interesse compreende vulnerabilidades classificáveis por CWE em projetos Python. O corpus experimental será a versão 1.0 do RealVuln descrita por \citeonline{Pellew2026RealVuln}, composta por 26 repositórios Python intencionalmente vulneráveis e 817 entradas, 697 vulnerabilidades, 120 armadilhas. Essa versão será usada integralmente, caracterizando um censo do \textit{release} escolhido, e não uma nova amostra construída pelo autor.

A escolha do RealVuln decorre de sua correspondência direta com o objetivo do trabalho: o benchmark foi projetado para comparar ferramentas baseadas em regras e scanners baseados em LLM sobre os mesmos repositórios, publica o código-fonte, o \textit{ground truth}, o algoritmo de correspondência e os scripts de pontuação, além de incluir exemplos seguros necessários ao cálculo de falsos positivos \cite{Pellew2026RealVuln}. O VUDENC permanece relevante como evidência histórica sobre detecção em Python, mas não será o corpus principal deste experimento porque sua rotulagem automática e sua granularidade em janelas de tokens exigiriam uma reconstrução de \textit{ground truth} incompatível com a comparação em nível de repositório.

A unidade de análise é a entrada rotulada do RealVuln: uma vulnerabilidade conhecida ou um falso-positivo armadilha, associada a repositório, arquivo, intervalo de linhas, CWE e evidência. A unidade de execução é um repositório completo analisado por uma condição experimental. A unidade de observação é o achado individual emitido pela condição, posteriormente associado, ou não, a uma entrada do \textit{ground truth}. Essa distinção permite que o agente navegue pelo projeto inteiro sem reduzir a avaliação a funções isoladas.

Será fixado o identificador da versão do benchmark, o \textit{commit} do repositório RealVuln, o \textit{commit} de cada projeto-alvo e o \textit{hash} dos arquivos de \textit{ground truth}. Nenhum repositório ou rótulo será removido após a observação dos resultados. Caso uma falha técnica impeça determinada condição de concluir um projeto, o modo estrito do benchmark será aplicado: as vulnerabilidades daquele projeto serão contabilizadas como não detectadas, e a falha de execução será reportada separadamente.

\section{Controle do Corpus e Cegamento}\label{sec:repositorio-experimental}

O RealVuln não será alterado nem terá suas vulnerabilidades recriadas. Para cada repositório, será utilizada a revisão indicada pelo manifesto oficial, e as seis condições receberão cópias com o mesmo conteúdo verificável por \textit{hash}. Código-fonte, configurações, dependências declaradas e testes pertencentes à revisão serão preservados; metadados de versionamento, artefatos gerados e documentos que revelem explicitamente as respostas serão excluídos por uma regra definida antes da coleta e aplicada de forma idêntica a todas as condições.

O \textit{ground truth}, o avaliador e a correspondência entre os projetos originais e seus identificadores opacos permanecerão fisicamente separados das áreas entregues às ferramentas. Cada execução utilizará uma área de trabalho nova. As condições de IA receberão apenas o código e a instrução padronizada; nas condições híbridas, será acrescentada somente a união normalizada dos alertas de C1 e C2, sem indicação de verdadeiro ou falso positivo. Esses controles reduzem a exposição direta das respostas e evitam estado residual entre condições e repetições.

Um manifesto interno registrará o alvo, a condição, a repetição, o \textit{commit}, o \textit{hash} da entrada e o estado da execução. Ele permitirá retomar a fila após interrupções sem reutilizar áreas de trabalho ou sessões anteriores. A estrutura de diretórios e os mecanismos utilizados para implementar essa separação são descritos no Capítulo~\ref{cap:desenvolvimento}, Seção~\ref{sec:organizacao-repositorio-desenvolvimento}.

\section{Importação e Uso do \textit{Ground Truth}}\label{sec:ground-truth-metodologia}

O \textit{ground truth} será importado diretamente da versão fixada do RealVuln, sem reconstrução ou alteração de rótulos pelo autor. Cada entrada contém identificador, indicação de vulnerabilidade ou armadilha segura, CWE principal, conjunto de CWEs aceitáveis, arquivo, intervalo de linhas, severidade e evidência de rotulagem \cite{Pellew2026RealVuln}. Antes da coleta, o arquivo será validado pelo verificador disponibilizado pelo benchmark e seu \textit{hash} será registrado.

Será preservado o algoritmo de correspondência do benchmark: caminho de arquivo normalizado, CWE pertencente ao conjunto aceitável da entrada e linha dentro da tolerância definida pelo RealVuln. Cada entrada poderá ser consumida uma única vez, evitando que alertas duplicados ampliem artificialmente a detecção. As saídas de Bandit, Semgrep, Cursor e Codex serão apenas convertidas para o esquema de entrada esperado pelo avaliador; a lógica de pontuação não será reimplementada quando o script oficial puder ser reutilizado.

Conforme as regras do RealVuln, um achado associado a uma entrada vulnerável será verdadeiro positivo; um achado associado a uma armadilha segura ou não associado a nenhuma entrada será falso positivo; uma vulnerabilidade sem correspondência será falso negativo; e uma armadilha corretamente ignorada será verdadeiro negativo. Eventuais erros identificados no \textit{ground truth} serão registrados e comunicados ao projeto, mas não corrigidos silenciosamente durante a coleta. Qualquer análise com rótulo corrigido será apresentada separadamente como análise de sensibilidade, preservando o resultado obtido com a versão oficial.

\section{Ambiente de Execução e Docker}\label{sec:materiais-ambiente}

A Tabela~\ref{tab:ambiente-experimental} apresenta os principais materiais e controles do ambiente. A justificativa técnica das ferramentas foi desenvolvida no Capítulo~\ref{cap:fundTeoric}; aqui são registrados apenas os elementos necessários à reprodução das execuções.

\begin{table}[htb]
\ABNTEXfontereduzida
\caption{\label{tab:ambiente-experimental}Materiais e controles do ambiente experimental}
\centering
\begin{tabularx}{\textwidth}{|l|X|X|}
\hline
\textbf{Categoria} & \textbf{Item} & \textbf{Controle registrado} \\
\hline
SAST & Bandit & Versão, comando, configuração, imagem Docker, duração e saída JSON \\
\hline
SAST & Semgrep Community Edition & Versão, comando, regras locais, \textit{hash} das regras, imagem Docker, duração e saída JSON \\
\hline
Conteinerização & Docker & Versão do motor, imagem-base por \textit{digest}, limites de recursos, rede e volumes \\
\hline
IA empresarial & Cursor & Versão do cliente, estratégia de seleção, modelo solicitado e efetivamente exibido, nível de raciocínio, modo de contexto, data e hora, instrução, permissões, duração, consumo de tokens informado e identificador da sessão \\
\hline
IA empresarial & OpenAI Codex & Versão do cliente, modelo solicitado e efetivamente exibido, nível de raciocínio, modo de contexto, data e hora, instrução, permissões, duração, consumo de tokens informado e identificador da sessão \\
\hline
Dados & RealVuln v1.0 & 26 repositórios, manifesto, 817 entradas rotuladas, \textit{ground truth} e avaliador fixados por \textit{hash} \\
\hline
Infraestrutura & \textit{Harness} em Python & Código versionado, esquema comum, registros brutos, normalização e cálculo das métricas \\
\hline
\end{tabularx}
\legend{Fonte: Elaborado pelo autor (2026).}
\end{table}

Bandit e Semgrep serão executados a partir de uma imagem Docker imutável, com versões fixadas, código-alvo somente para leitura e saída gravada separadamente. Cada combinação entre ferramenta e repositório utilizará um contêiner novo, sem rede, com usuário sem privilégios e os mesmos limites de CPU, memória e tempo. Assim, a conteinerização atua como variável de controle das condições SAST, enquanto o manifesto externo ao contêiner permite retomar a fila de execuções.

Em conformidade com a delimitação estabelecida na Seção~\ref{subsec:ferramentas-sast}, não serão criadas regras específicas para os repositórios do corpus. O Bandit será executado com seu conjunto nativo de \textit{plugins} e sem supressões ajustadas após a observação dos resultados. O Semgrep Community Edition utilizará um conjunto público de regras de segurança para Python, selecionado antes do estudo-piloto, copiado para \texttt{docker/regras-semgrep/} e fixado por versão e \textit{hash}. O nome do conjunto, os arquivos efetivamente carregados e a linha de comando completa serão publicados, pois a expressão ``configuração padrão'' não identifica sozinha quais regras públicas foram aplicadas pelo Semgrep.

A conteinerização equaliza o ambiente local das condições C1 e C2, mas não torna os serviços de IA equivalentes nem imutáveis. Por isso, cada bloco de C3--C6 será executado em uma janela de coleta tão curta quanto possível, com sessões novas, o mesmo texto de instrução, a mesma cópia do repositório e permissões equivalentes. No bloco principal, o modelo será fixado em GPT-5.6 Luna nos dois produtos. No bloco complementar, se realizado, Cursor Auto e Codex com GPT-5.6 Terra serão analisados como configurações próprias de cada produto, sem atribuir suas diferenças exclusivamente à camada de ferramenta. Alterações de modelo ou de produto durante a coleta interromperão o bloco afetado, que será reiniciado ou analisado separadamente. A construção da imagem, os comandos e os adaptadores pertencem à implementação e são detalhados no Capítulo~\ref{cap:desenvolvimento}.

\section{Instrumentação e Dados Coletados}\label{subsec:implementacao-harness}

Um \textit{harness} em Python será utilizado para preparar entradas idênticas, controlar a fila, preservar saídas brutas e metadados e converter os achados para o formato aceito pelo avaliador. Para cada execução serão registrados, no mínimo, condição, ferramenta, bloco de modelos, alvo, repetição, arquivo, intervalo de linhas, regra ou CWE informado, severidade, descrição, recomendação, duração e estado da execução. Nas condições baseadas em IA, serão registrados também a estratégia de seleção, o modelo solicitado, o modelo efetivamente exibido e os consumos de tokens de entrada, de saída e total quando esses dados forem disponibilizados pelo produto.

As condições C3--C6 deverão emitir uma lista JSON com caminho do arquivo, linha inicial e final, CWE, severidade, descrição da causa e recomendação. A lista vazia será aceita como ausência de achados. Campos omitidos permanecerão nulos: o normalizador não utilizará o \textit{ground truth} para completar localizações, CWE ou severidades. A união SAST fornecida às condições híbridas será deduplicada antes da nova execução dos agentes, mantendo a origem dos alertas e sem anexar rótulos do benchmark.

Os adaptadores serão verificados antes da coleta com saídas artificiais externas ao corpus e com uma execução de fumaça descartável. A arquitetura do \textit{harness}, o esquema comum, as regras de deduplicação, os comandos e os testes de integração são especificados no Capítulo~\ref{cap:desenvolvimento}, Seções~\ref{sec:arquitetura-harness} e~\ref{sec:verificacao-artefato-desenvolvimento}.

\section{Métricas de Avaliação}\label{sec:metricas-metodologia}

\subsection{Métricas quantitativas de detecção}\label{subsec:metricas-quantitativas-metodologia}

A matriz de confusão será produzida pelo avaliador do RealVuln no nível das entradas rotuladas. Vulnerabilidades correspondidas serão verdadeiros positivos (TP), vulnerabilidades não encontradas serão falsos negativos (FN), armadilhas ou regiões não rotuladas sinalizadas serão falsos positivos (FP), e armadilhas corretamente ignoradas serão verdadeiros negativos (TN). A pontuação principal utilizará o modo estrito e agregado por microcontagem do benchmark, no qual uma falha completa de análise não é removida do denominador \cite{Pellew2026RealVuln}.

Serão calculadas, por condição, repositório, categoria CWE e severidade:

\begin{itemize}
\item \textbf{Precisão}: $TP/(TP+FP)$;
\item \textbf{Recall}: $TP/(TP+FN)$;
\item \textbf{F1}: média harmônica entre precisão e \textit{recall};
\item \textbf{F2}: média harmônica ponderada com maior peso para o \textit{recall}, adequada à maior criticidade dos falsos negativos em segurança \cite{Kalouptsoglou2023VulnPrediction};
\item \textbf{F3}: média harmônica ponderada que atribui peso nove vezes maior ao \textit{recall}, métrica principal do RealVuln \cite{Pellew2026RealVuln};
\item \textbf{Taxa de falsos positivos}: $FP/(FP+TN)$.
\end{itemize}

Além dos valores pontuais, serão apresentados intervalos de confiança de 95\% obtidos por reamostragem no nível do repositório, preservando o agrupamento de entradas pertencentes ao mesmo projeto. As diferenças entre condições aplicadas aos mesmos repositórios serão calculadas de forma pareada, com intervalos obtidos pelo mesmo procedimento de reamostragem. A análise por CWE e severidade será apresentada mesmo quando a quantidade de entradas permitir apenas interpretação descritiva.

\subsection{Tempo de execução e consumo de tokens}\label{subsec:tempo-tokens-metodologia}

O tempo de execução será medido pelo \textit{harness} como tempo de relógio entre o envio da tarefa e a disponibilidade da saída final. Nas condições SAST, serão registrados também o tempo de inicialização do contêiner e o tempo interno do analisador, permitindo distinguir a sobrecarga da infraestrutura do tempo da análise. Nas condições de IA, o tempo incluirá as iterações autônomas do agente até sua resposta final ou até o limite de execução previamente definido. Serão reportados mediana, intervalo interquartil e valores mínimos e máximos por condição, pois tempos de serviços remotos podem apresentar distribuição assimétrica.

O consumo de tokens será coletado apenas nas condições baseadas em IA e conforme os contadores efetivamente disponibilizados por cada produto. Serão preservados separadamente tokens de entrada, tokens de saída e total de tokens; categorias adicionais expostas pela ferramenta, como tokens em cache ou de raciocínio, serão mantidas em campos próprios. Quando um produto não informar algum desses valores, o campo será registrado como indisponível, sem estimativa a partir de caracteres, palavras ou duração. Os resultados serão resumidos por condição e por repositório, com mediana, intervalo interquartil e valores mínimos e máximos. Como produtos distintos podem empregar tokenizadores e mecanismos de contabilização diferentes, a comparação será descritiva e explicitará a disponibilidade e a definição dos contadores usados.

\subsection{Repetição e estabilidade}\label{subsec:testes-ensaios}

As condições determinísticas C1 e C2 serão executadas uma vez por repositório, após a validação da imagem Docker. As condições baseadas em IA, C3--C6, serão executadas três vezes por repositório em cada bloco efetivamente realizado, sempre em sessão nova e sem memória das execuções anteriores. A ordem dos repositórios será aleatorizada por bloco, e todas as execuções de um bloco utilizarão as mesmas versões e configurações observáveis. As repetições do bloco complementar não serão combinadas com as do bloco principal para calcular estabilidade.

As métricas quantitativas serão calculadas separadamente em cada repetição e resumidas por média, amplitude e desvio entre execuções. No nível de cada entrada vulnerável, será reportada a frequência de detecção em zero, uma, duas ou três execuções. Essa frequência representa estabilidade empírica e não será convertida silenciosamente em um único resultado por voto majoritário. Para tornar a avaliação qualitativa exequível e evitar seleção retrospectiva da melhor resposta, a primeira execução, definida antes da coleta, será a saída principal da rubrica; as demais serão usadas na análise de estabilidade.

\subsection{Avaliação qualitativa dos relatórios}\label{subsec:metricas-qualitativas-metodologia}

A avaliação qualitativa será aplicada aos relatórios verdadeiros-positivos da execução principal. Para manter a etapa exequível, primeiro serão sorteadas até 20 entradas detectadas por todas as condições, que formarão o recorte diretamente comparável. Em seguida, a amostra de cada condição será completada até 50 relatórios com sorteio estratificado por CWE e severidade. Quando uma condição produzir menos de 50 verdadeiros-positivos, todos serão avaliados. A semente e a lista sorteada serão publicadas. Assim, serão apresentados tanto um recorte comum, sobre as mesmas vulnerabilidades, quanto um recorte representativo da utilidade efetivamente entregue por cada condição. O sorteio e a aplicação da rubrica serão realizados separadamente para cada bloco de modelos, sem reunir respostas Luna, Auto e Terra em uma única amostra.

A Tabela~\ref{tab:dimensoes-rubrica} define as três dimensões ordinais, pontuadas de 1 a 5. A clareza será avaliada separadamente por checklist estrutural.

\begin{table}[htb]
\ABNTEXfontereduzida
\caption{\label{tab:dimensoes-rubrica}Dimensões da rubrica qualitativa}
\centering
\begin{tabularx}{\textwidth}{|l|X|X|}
\hline
\textbf{Dimensão} & \textbf{Critério} & \textbf{Fundamentação} \\
\hline
Correção técnica & Explica corretamente a causa e o fluxo da vulnerabilidade, sem afirmar relações inexistentes no código & \cite{Mao2025ExplainableVD} \\
\hline
Localização & Identifica o arquivo e o trecho responsável com granularidade compatível com o \textit{ground truth} & \cite{Zhang2023LocateDescribe, Mao2025ExplainableVD} \\
\hline
Acionabilidade & Propõe correção tecnicamente válida, específica e aplicável ao contexto analisado & \cite{Mao2025ExplainableVD, Aloraini2019SASTWarnings} \\
\hline
\end{tabularx}
\legend{Fonte: Elaborado pelo autor com base nas referências indicadas (2026).}
\end{table}

Os descritores-âncora da escala são apresentados na Tabela~\ref{tab:ancoras-rubrica}. Os níveis 2 e 4 representam posições intermediárias entre os descritores adjacentes.

\begin{table}[htb]
\ABNTEXfontereduzida
\caption{\label{tab:ancoras-rubrica}Descritores-âncora da rubrica qualitativa}
\centering
\begin{tabularx}{\textwidth}{|c|X|X|X|}
\hline
\textbf{Nível} & \textbf{Correção técnica} & \textbf{Localização} & \textbf{Acionabilidade} \\
\hline
1 & Erro factual que compromete a explicação da causa & Localização ausente ou incorreta & Ausência de recomendação ou proposta incompatível com o caso \\
\hline
3 & Causa geral correta, com imprecisões secundárias & Arquivo ou função corretos, sem o trecho exato & Recomendação válida, porém genérica \\
\hline
5 & Causa e fluxo tecnicamente corretos e completos & Arquivo e trecho vulnerável corretamente identificados & Correção específica, tecnicamente válida e diretamente aplicável \\
\hline
\end{tabularx}
\legend{Fonte: Elaborado pelo autor (2026).}
\end{table}

\subsubsection{Operacionalização da clareza}\label{subsubsec:operacionalizacao-clareza}

A clareza será medida pelo checklist da Tabela~\ref{tab:checklist-clareza}. Embora as respostas sejam binárias, os critérios ainda envolvem julgamento humano; por isso, o escore de zero a cinco será tratado como contagem estruturada, e não como medida inteiramente objetiva.

\begin{table}[htb]
\ABNTEXfontereduzida
\caption{\label{tab:checklist-clareza}Checklist estrutural de clareza dos relatórios}
\centering
\begin{tabularx}{\textwidth}{|l|X|}
\hline
\textbf{Critério} & \textbf{Pergunta binária} \\
\hline
C1 --- Problema & O texto declara qual é a vulnerabilidade, além de apresentar apenas regra ou CWE? \\
\hline
C2 --- Risco & O texto explica impacto ou cenário de exploração associado à falha? \\
\hline
C3 --- Contexto técnico & O jargão empregado é compreensível no contexto do relatório? \\
\hline
C4 --- Concretude & O texto menciona elementos concretos do código, como função, variável ou parâmetro? \\
\hline
C5 --- Concisão & O texto comunica problema, risco e contexto sem omissões impeditivas nem repetição desnecessária? \\
\hline
\end{tabularx}
\legend{Fonte: Elaborado pelo autor (2026).}
\end{table}

Como medida secundária, será calculado o Índice de Legibilidade de Flesch adaptado ao português brasileiro \cite{Martins1996ReadabilityPortuguese}. O idioma das respostas será fixado pelo \textit{prompt}; trechos de código, caminhos e identificadores serão removidos do cálculo de legibilidade. O índice será interpretado apenas como facilidade sintática de leitura e não como substituto da avaliação técnica ou do checklist.

Os relatórios serão apresentados aos avaliadores sem identificação direta da ferramenta e em ordem aleatorizada, sempre que a remoção da origem não eliminar conteúdo necessário à avaliação. O autor avaliará a totalidade do recorte qualitativo, e um segundo avaliador examinará independentemente pelo menos 20\% dos relatórios, em amostra estratificada por condição e CWE. Para as dimensões ordinais será reportado o Kappa ponderado de Cohen; para os itens binários de clareza, percentual de concordância e Kappa não ponderado \cite{Landis1977ObserverAgreement}. Divergências não serão resolvidas alterando retroativamente os descritores, mas serão registradas e discutidas.

\subsection{Avaliação de severidade e priorização}\label{subsec:priorizacao-metodologia}

A presença de uma indicação de severidade será reportada para todas as condições. A correção da priorização será comparada com a categoria de severidade publicada em cada entrada vulnerável do RealVuln. As categorias produzidas pelas ferramentas serão normalizadas para a mesma escala ordinal do benchmark (baixa, média, alta e crítica), mantendo-se também o valor bruto fornecido pela ferramenta.

Serão calculadas a cobertura de priorização, a concordância exata com a categoria de referência, a concordância com tolerância de uma categoria e o Kappa ponderado. Achados sem severidade receberão valor ausente e reduzirão a cobertura, mas não terão sua categoria inferida pelo \textit{harness}. Essa delimitação impede que a capacidade de exibir uma etiqueta seja confundida com a capacidade de priorizar corretamente.

\section{Protocolo Experimental}\label{sec:protocolo-experimental}

O protocolo será executado na seguinte ordem:

\begin{enumerate}
\item Fixar o \textit{commit} da versão 1.0 do RealVuln e registrar os \textit{hashes} do manifesto, do \textit{ground truth} e do avaliador oficial.
\item Executar o validador do benchmark e baixar os 26 repositórios nos \textit{commits} indicados pelo manifesto.
\item Gerar cópias sanitizadas com identificadores opacos e verificar que as áreas de trabalho não expõem o diretório \texttt{ground-truth}, o avaliador ou os manifestos internos.
\item Construir a imagem Docker, fixar versões e regras, e validar os adaptadores com saídas artificiais externas ao corpus.
\item Confirmar no estudo-piloto a disponibilidade do GPT-5.6 Luna nos dois produtos, fixar suas opções e decidir, antes de examinar os resultados principais, se as cotas e o limite de tokens permitem executar integralmente o bloco complementar.
\item Criar a fila completa de combinações entre condição, repositório e repetição, com ordem aleatorizada dentro de cada bloco.
\item Executar C1 e C2 em contêineres efêmeros, preservando comandos, duração, metadados e JSON bruto.
\item Normalizar e deduplicar os alertas SAST, sem consultar os resultados das condições de IA.
\item Executar C3 e C4 em três sessões independentes por repositório, sem acesso aos alertas SAST ou ao \textit{ground truth}.
\item Executar C5 e C6 em três sessões independentes por repositório, fornecendo apenas os alertas SAST normalizados e o código correspondente.
\item Se aprovado no estudo-piloto, repetir C3--C6 em bloco complementar completo, usando Cursor Auto e GPT-5.6 Terra no Codex, sem alterar as demais instruções ou permissões.
\item Converter as seis condições para o esquema comum e executar o avaliador oficial do RealVuln no modo estrito.
\item Calcular métricas de detecção, estabilidade, tempo, consumo de tokens, intervalos de confiança e comparações pareadas.
\item Sortear a amostra qualitativa da primeira execução, aplicar a rubrica e calcular a concordância entre avaliadores.
\item Reportar resultados por condição, repositório, CWE e severidade, preservando tanto os dados agregados quanto as falhas de execução.
\end{enumerate}

Prompts, configurações, imagem Docker, regras do Semgrep, versão do Bandit, scripts de normalização, manifestos permitidos pela licença dos dados e resultados brutos serão versionados como artefatos de reprodução do estudo. O \texttt{oracle/} será preservado para auditoria, mas permanecerá separado das áreas de trabalho usadas durante a coleta.

\section{Viabilidade da Pesquisa}\label{sec:viabilidade-pesquisa}

A adoção integral de uma versão publicada do RealVuln elimina a construção manual do corpus e do \textit{ground truth}. Os manifestos permitem baixar os projetos uma vez, verificar sua integridade e retomar execuções interrompidas sem reconstruir cada alvo. A conteinerização reduz o esforço de reinstalação e a variação do ambiente SAST. A viabilidade das condições de IA depende de licenças, cotas e limites de uso existentes no período da coleta; por serem características mutáveis dos produtos, elas serão verificadas e registradas no estudo-piloto, sem pressupor disponibilidade ilimitada. O bloco Luna constitui o núcleo obrigatório; a expansão para Cursor Auto e Codex com GPT-5.6 Terra somente será incorporada se o limite de tokens reservado comportar sua execução integral.

Os principais fatores de esforço passam a ser a integração dos formatos de saída, as três repetições das condições de IA e a avaliação qualitativa. O corpus quantitativo permanecerá completo; para controlar o esforço humano, somente a avaliação qualitativa utilizará a amostra estratificada definida na Seção~\ref{subsec:metricas-qualitativas-metodologia}.

\section{Validade e Limitações}\label{sec:validade-limitacoes}

A Tabela~\ref{tab:ameacas-validade} sintetiza as principais ameaças e os controles previstos.

\begin{table}[htb]
\ABNTEXfontereduzida
\caption{\label{tab:ameacas-validade}Ameaças à validade e estratégias de mitigação}
\centering
\begin{tabularx}{\textwidth}{|X|X|}
\hline
\textbf{Ameaça} & \textbf{Mitigação} \\
\hline
Variação do ambiente SAST & Imagem Docker imutável, regras locais, código somente leitura e registro de comandos e \textit{hashes} \\
\hline
Estado residual entre repositórios & Área de trabalho, contêiner e sessão novos para cada execução, com retomada controlada por manifesto \\
\hline
Exposição direta da resposta esperada & Identificadores opacos e separação física do \textit{ground truth} e do avaliador \\
\hline
Mudança de modelo ou produto de IA & Coleta em blocos curtos, registro de versão/modelo/data e reinício ou separação de blocos afetados \\
\hline
Confundimento entre ferramenta e modelo & GPT-5.6 Luna fixado nos dois produtos no bloco principal; configurações nativas examinadas apenas em bloco complementar separado \\
\hline
Uso inadvertido de SAST pelo braço IA & Restrições explícitas de ferramentas, inspeção dos registros de ações e descarte documentado de execução contaminada \\
\hline
Não determinismo & Três sessões independentes, ordem aleatorizada e reporte da frequência de detecção por entrada \\
\hline
Rótulos ruidosos ou incompletos & Versão e \textit{hash} fixados, validador oficial e análise de sensibilidade para correções documentadas \\
\hline
Memorização do benchmark por modelos & Código opaco no ambiente de execução e reconhecimento explícito de que o isolamento não elimina possível exposição durante o treinamento \\
\hline
Viés na adjudicação e na rubrica & Cegamento da ferramenta, descritores prévios, segundo avaliador e reporte da concordância \\
\hline
Granularidades diferentes de alerta & Adaptadores para o esquema comum e reutilização do algoritmo de correspondência do RealVuln \\
\hline
Alta densidade de vulnerabilidades & Resultados interpretados como desempenho no benchmark, sem inferir a prevalência ou a precisão esperada em projetos de produção \\
\hline
Generalização limitada & Conclusões restritas a Python, ao RealVuln v1.0, às versões, regras, prompts e ferramentas efetivamente avaliados \\
\hline
\end{tabularx}
\legend{Fonte: Elaborado pelo autor (2026).}
\end{table}

A validade interna é fortalecida pelo uso dos mesmos repositórios e do mesmo avaliador em todas as condições, pelo controle do ambiente SAST, pelo modelo comum no bloco principal e pela preservação das saídas brutas. A validade de construto depende de não reduzir ``qualidade'' a uma única métrica: por isso, detecção, estabilidade, explicação e priorização são avaliadas separadamente, enquanto tempo e consumo de tokens caracterizam operacionalmente as execuções. A validade da medida de tokens depende dos contadores expostos pelos produtos e de suas definições, que podem não ser diretamente equivalentes. No bloco complementar, a diferença entre Cursor Auto e Codex com GPT-5.6 Terra representa a configuração conjunta produto--modelo e não permite isolar causalmente a camada da ferramenta. A validade externa permanece limitada por o RealVuln v1.0 ser composto por aplicações educacionais e CTF deliberadamente vulneráveis, com densidade superior à de projetos de produção, além da rápida evolução dos produtos comerciais. Consequentemente, o estudo sustentará conclusões sobre o benchmark e as configurações observadas durante a coleta, e não sobre uma capacidade permanente ou universal de Cursor, Codex, Bandit, Semgrep ou de todas as abordagens híbridas.

%\chapter{Desenvolvimento}\label{cap:desenvolvimento}
\addcontentsline{toc}{chapter}{Desenvolvimento}

Este capítulo descreve a implementação efetivamente utilizada para executar o protocolo definido no Capítulo~\ref{cap:metodologia}. O artefato resultante é um repositório controlador que adquire e verifica o corpus, prepara entradas sanitizadas, executa as seis condições experimentais, valida as respostas, preserva todas as tentativas e produz evidências auditáveis. Diferentemente de uma especificação anterior à coleta, o texto a seguir registra as versões, os controles, os formatos e os desvios observados durante a execução.

A coleta das seis condições foi concluída sobre os 26 repositórios do RealVuln v1.0. Foram selecionadas 364 unidades válidas: 52 execuções SAST, correspondentes a uma execução de C1 e C2 por alvo, e 312 execuções de IA, correspondentes a três repetições de C3--C6 por alvo. As retentativas não foram sobrescritas; somadas às unidades selecionadas, produziram 433 manifestos ou tentativas preservadas. Esses números caracterizam a completude operacional do artefato. As medidas de detecção e a comparação com o \textit{ground truth} pertencem ao Capítulo~\ref{cap:avaliacao}.

\section{Visão Geral da Implementação}\label{sec:visao-geral-desenvolvimento}

O controlador foi implementado predominantemente em Python e PowerShell. Python concentra os modelos de dados, a validação dos contratos, a preparação do corpus, a fila, a execução e a normalização de C1 e C2. PowerShell coordena a aquisição, o Docker Desktop e as interfaces de linha de comando do Cursor e do Codex. Os projetos vulneráveis foram tratados como dados: não foram importados, testados nem executados pelo controlador.

O fluxo foi dividido em três domínios. O domínio de entrada contém as árvores sanitizadas e os \textit{payloads} de código entregues às ferramentas. O domínio de controle contém configurações congeladas, filas, \textit{prompts}, esquemas, adaptadores e áreas de trabalho. O domínio de avaliação contém o \textit{ground truth} e o validador oficial do RealVuln. Durante a coleta, os executores de C1--C6 não consultaram o domínio de avaliação; as auditorias das condições de IA registraram explicitamente \texttt{oracle\_consultado=false}. Essa separação materializa o cegamento descrito na Seção~\ref{sec:repositorio-experimental}.

A cadeia implementada pode ser resumida em: aquisição das revisões fixadas; sanitização e inventário; congelamento das configurações; execução serial e preservação das tentativas; validação sintática e estrutural; auditoria de fechamento; e, em etapa separada, pontuação contra o \textit{oracle}. A separação entre coleta e pontuação impede que rótulos do benchmark influenciem a geração ou a normalização dos achados.

\section{Organização do Repositório Controlador}\label{sec:organizacao-repositorio-desenvolvimento}

A Tabela~\ref{tab:estrutura-repositorio-experimental} apresenta a estrutura efetivamente adotada. Três diretórios possuem papéis distintos na cadeia de custódia: \texttt{execucoes/} mantém o estado e as áreas mínimas de trabalho; \texttt{resultados/} contém os registros primários por tentativa; e \texttt{evidencias/} contém sínteses derivadas e inventários ligados aos registros primários por SHA-256.

\begin{table}[htb]
\ABNTEXfontereduzida
\caption{Estrutura efetiva do repositório controlador}
\label{tab:estrutura-repositorio-experimental}
\centering
\begin{tabularx}{\textwidth}{|l|X|X|}
\hline
\textbf{Diretório} & \textbf{Conteúdo} & \textbf{Regra de acesso e uso} \\
\hline
\texttt{config/} & \textit{Locks} de fontes, corpus, fila, sanitização, agentes, \textit{prompt} e esquema de saída & Contratos congelados; alterações posteriores invalidam a comparação com a coleta realizada \\
\hline
\texttt{benchmark/} & \textit{Cache} local reconstruível do RealVuln e dos repositórios de origem & Usado somente na aquisição; não é entregue às ferramentas nem redistribuído \\
\hline
\texttt{alvos/} & Árvores sanitizadas identificadas como \texttt{ALVO-NNNN} & Tratadas como dados e verificadas por inventário antes do uso \\
\hline
\texttt{oracle/} & Manifesto, \textit{ground truth} e validador oficiais & Segregado da coleta; acesso reservado à validação do benchmark e à pontuação posterior \\
\hline
\texttt{runner/} e \texttt{scripts/} & Núcleo Python, adaptadores e orquestradores PowerShell & Executam e validam o protocolo sem completar campos com informação do \textit{oracle} \\
\hline
\texttt{docker/} & Imagem, dependências fixadas e cópia local das regras Semgrep & Usado por C1 e C2; o pacote de regras é reconstruível, mas não é redistribuído \\
\hline
\texttt{execucoes/} & Fila, alertas híbridos congelados e áreas de trabalho das chamadas & Estado mutável do controlador; não é fonte única de métricas \\
\hline
\texttt{resultados/} & Saídas brutas, normalizadas, eventos, respostas e manifestos por tentativa & Registro primário e imutável; nenhuma tentativa é reutilizada ou sobrescrita \\
\hline
\texttt{evidencias/} & Inventários, auditorias de fechamento e sínteses derivadas & Evidência publicada, vinculada aos primários por hashes \\
\hline
\end{tabularx}
\legend{Fonte: Elaborado pelo autor (2026).}
\end{table}

Em C1 e C2, cada tentativa foi gravada em \path{resultados/CN-ALVO-NNNN-R01/tentativa-NNN/}. Em C3--C6, o caminho equivalente foi dividido por finalidade, por exemplo \path{resultados/ia/coleta/CN-ALVO-NNNN-RNN/tentativa-NNN/}. A finalidade \texttt{piloto} permaneceu em um \textit{namespace} separado e, por isso, não pôde ocupar uma unidade da coleta principal.

A publicação do artefato considerou também as licenças e o conteúdo sensível dos projetos de terceiros. Dezesseis árvores sanitizadas, com licença permissiva ou copyleft identificada, foram incluídas integralmente em \texttt{alvos/corpus-v1/}. Os outros dez alvos permaneceram reconstruíveis por URL, \textit{commit}, tamanho e SHA-256 de cada arquivo. Nove não apresentavam licença explícita na revisão adquirida e um continha uma chave privada de demonstração. O inventário publicado cobre ainda \texttt{oracle/}, \texttt{execucoes/} e \texttt{resultados/}, permitindo verificar a integridade do conjunto redistribuído sem publicar material não autorizado.

\section{Preparação do Corpus Experimental}\label{sec:preparacao-corpus}

\subsection{Fixação do RealVuln e verificação de integridade}\label{subsec:fixacao-realvuln}

O RealVuln foi fixado na tag \texttt{v1.0}, cujo objeto anotado é identificado pelo prefixo \texttt{aa9f7321c8c5} e cujo \textit{commit} resolvido possui o prefixo \texttt{d98e9fc91273}. O arquivo \texttt{config/corpus-realvuln-v1.lock.json} registra a URL e o \textit{commit} completo dos 26 projetos. Três URLs de origem estavam indisponíveis durante a aquisição; nesses casos foram usados espelhos previamente autorizados apenas como transporte, com aceitação condicionada à presença do mesmo objeto Git fixado no manifesto.

O script de aquisição obteve cada revisão sem executar \textit{hooks}, submódulos, testes ou arquivos do alvo. Antes da promoção das áreas temporárias, foram conferidos o manifesto, os \textit{commits}, os arquivos regulares permitidos e os hashes da árvore produzida. O validador oficial \texttt{validate\_gt.py} terminou com código zero sobre os 26 arquivos de \textit{ground truth}, totalizando 817 entradas: 697 vulnerabilidades e 120 armadilhas de falso positivo. O artefato preserva tanto essa saída literal quanto uma auditoria por repositório.

\subsection{Sanitização e identificadores opacos}\label{subsec:sanitizacao-alvos}

A política \texttt{sanitizacao-v1} removeu metadados \texttt{.git}, ambientes virtuais, dependências instaladas, \textit{caches}, artefatos de compilação, \textit{ground truth}, \textit{oracle}, soluções, \textit{walkthroughs} e \textit{write-ups}. Apenas objetos Git com modo de arquivo regular foram aceitos; três ligações simbólicas conhecidas foram omitidas sem que seus destinos fossem seguidos. A política foi congelada em \texttt{config/politica-sanitizacao-v1.json}.

Cada projeto recebeu um identificador opaco entre \texttt{ALVO-0001} e \texttt{ALVO-0026}. A correspondência entre esses identificadores e os projetos originais permaneceu no \textit{lock} do corpus, fora das áreas de trabalho dos agentes. Para cada alvo, o preparador produziu uma lista ordenada contendo caminho relativo, tamanho e SHA-256 de cada arquivo, além de um SHA-256 canônico da árvore. O executor verificou esse identificador antes e depois de C1 e C2; os executores de IA reconstruíram o \textit{payload} apenas com arquivos Python presentes no inventário congelado.

Essa estratégia também evitou depender de data de modificação ou de outros metadados do sistema de arquivos. A igualdade passou a ser definida pelos bytes e pelos caminhos relativos. As evidências em \texttt{evidencias/corpus-realvuln-v1/} permitem reconstruir e comparar todos os 26 alvos, inclusive aqueles que não puderam ser redistribuídos integralmente.

\section{Arquitetura do \textit{Harness}}\label{sec:arquitetura-harness}

\subsection{Módulos e responsabilidades}\label{subsec:modulos-harness}

A Tabela~\ref{tab:modulos-harness} relaciona os módulos implementados às suas responsabilidades. A decomposição separa aquisição, execução e avaliação, de modo que o código que entrega entradas às ferramentas não precise consultar o \textit{ground truth}.

\begin{table}[htb]
\ABNTEXfontereduzida
\caption{Módulos implementados no \textit{harness}}
\label{tab:modulos-harness}
\centering
\begin{tabularx}{\textwidth}{|l|X|X|}
\hline
\textbf{Componente} & \textbf{Responsabilidade implementada} & \textbf{Saída principal} \\
\hline
\texttt{aquisicao.py} e \texttt{corpus.py} & Validar fontes, revisões, políticas, espelhos e \textit{ground truth} & \textit{Locks}, auditorias e manifesto do corpus \\
\hline
\texttt{preparacao.py} & Copiar apenas arquivos permitidos e calcular inventários canônicos & Árvore sanitizada e inventário SHA-256 \\
\hline
\texttt{fila.py} e \texttt{coleta.py} & Controlar estados, exclusão mútua, retomada e fechamento de C1/C2 & Fila e histórico por tentativa \\
\hline
\texttt{executor\_sast.py} & Executar o scanner, medir duração, preservar bytes e fechar manifestos & Bruto, normalizado e manifesto terminal \\
\hline
\texttt{adaptadores/} & Converter Bandit e Semgrep para o esquema comum sem consultar o \textit{oracle} & Lista de achados normalizados \\
\hline
\texttt{scripts/ia/} & Preparar \textit{payloads}, chamar Cursor ou Codex e validar eventos e respostas & Resposta, eventos, metadados e manifesto de IA \\
\hline
\texttt{scripts/*-lote.ps1} & Planejar, executar serialmente, pular unidades válidas e retomar falhas & Resumo de lote e novas tentativas \\
\hline
\end{tabularx}
\legend{Fonte: Elaborado pelo autor (2026).}
\end{table}

Cada unidade recebeu um identificador formado por condição, alvo e repetição, como \texttt{C4-ALVO-0001-R02}. Em C1 e C2, a fila persistente admite os estados \texttt{pendente}, \texttt{em\_execucao}, \texttt{concluida} e \texttt{falha}, com histórico de transições e \textit{lock} de exclusão mútua. Em C3--C6, o lote revalida a tentativa existente antes de considerá-la concluída. Uma falha cria a próxima pasta \texttt{tentativa-NNN}; a pasta anterior permanece intacta.

Os registros terminais somente foram promovidos depois da gravação dos artefatos obrigatórios. Arquivos em produção receberam o sufixo \texttt{.part} e foram renomeados após o encerramento, enquanto JSONs de controle foram gravados de forma atômica. Os manifestos incluem hashes dos artefatos, horários, duração monotônica ou duração informada pelo adaptador, código de saída, versão observada, entrada e configuração. Esse desenho permitiu retomar as coletas sem confundir falha técnica com relatório vazio.

\subsection{Esquemas e preservação dos dados}\label{subsec:esquema-comum-achados}

As saídas nativas de Bandit e Semgrep foram preservadas primeiro em \texttt{bruto.json}. Somente depois o adaptador gerou \texttt{normalizado.json}. A Tabela~\ref{tab:esquema-achado-comum} resume o esquema dos achados SAST. Valores originais e normalizados coexistem; quando não existe mapeamento fechado, o campo normalizado permanece nulo.

\begin{table}[htb]
\ABNTEXfontereduzida
\caption{Esquema comum dos achados SAST normalizados}
\label{tab:esquema-achado-comum}
\centering
\begin{tabularx}{\textwidth}{|l|X|X|}
\hline
\textbf{Grupo} & \textbf{Campos principais} & \textbf{Tratamento} \\
\hline
Proveniência & condição, ferramenta, alvo, repetição, execução e hash da saída bruta & Obrigatórios e vinculados ao manifesto \\
\hline
Localização & caminho relativo, linha inicial e linha final & Normalização sintática; ausências não são inferidas \\
\hline
Classificação & regra, CWE, severidade e confiança, nas formas original e normalizada & Conversão apenas por mapas fechados \\
\hline
Relatório & descrição, evidência, recomendação e texto original & Conteúdo preservado sem complementação pelo controlador \\
\hline
Operação & início, término, duração, código de saída, estado e artefatos & Registrado no manifesto da tentativa \\
\hline
\end{tabularx}
\legend{Fonte: Elaborado pelo autor (2026).}
\end{table}

As respostas de C3--C6 utilizaram um contrato próprio, comum aos dois produtos. A raiz deve ser um \textit{array} JSON e cada item contém exatamente \texttt{arquivo}, \texttt{linha\_inicial}, \texttt{linha\_final}, \texttt{cwe}, \texttt{severidade}, \texttt{descricao} e \texttt{recomendacao}. Propriedades adicionais, caminhos ausentes do \textit{payload}, intervalos fora do arquivo, CWE malformada e severidade fora da enumeração invalidam a tentativa. A resposta bruta, contudo, permanece preservada mesmo quando a validação falha.

\section{Execução das Condições SAST (C1 e C2)}\label{sec:execucao-sast-desenvolvimento}

\subsection{Imagem Docker e isolamento}\label{subsec:construcao-imagem-docker}

C1 e C2 foram executadas com uma imagem experimental para \texttt{linux/amd64}, baseada em Python 3.12.13 sobre Debian Bookworm e fixada por \textit{digest}. A imagem continha Bandit 1.9.4 e Semgrep 1.172.0. As 69 dependências Python foram instaladas exclusivamente a partir de \textit{wheels} com versão e SHA-256 registrados. O corpus, os resultados e o \textit{oracle} não foram copiados para a imagem.

Cada tentativa utilizou um contêiner exclusivo com rede desativada, sistema de arquivos raiz somente leitura, usuário \texttt{10001:10001}, todas as \textit{capabilities} removidas, \texttt{no-new-privileges}, 2 CPUs, 3~GB de memória, 256 processos e \textit{tmpfs} de 64~MB com \texttt{noexec}, \texttt{nosuid} e \texttt{nodev}. A entrada foi montada como somente leitura e somente a pasta de saída recebeu escrita. O \texttt{docker.sock} não foi montado. O limite de cada scanner foi fixado em 900 segundos e as execuções ocorreram serialmente.

\subsection{Adaptador do Bandit}\label{subsec:execucao-bandit}

O Bandit 1.9.4 foi invocado recursivamente sobre \texttt{/entrada}, sem filtro posterior de severidade ou confiança:

\begin{verbatim}
bandit --recursive /entrada --format json \
  --output /saida/bruto.json.part
\end{verbatim}

Os códigos zero e um foram aceitos, pois o Bandit usa o código um para indicar que encontrou problemas. O adaptador preservou identificador do teste, mensagem, caminho, intervalo, severidade, confiança e CWE informado. A validação rejeitou JSON duplicado, valores não finitos e estruturas incompatíveis, sem alterar a saída original.

\subsection{Adaptador do Semgrep Community Edition}\label{subsec:execucao-semgrep}

O Semgrep Community Edition 1.172.0 utilizou apenas o diretório \texttt{python/} do repositório de regras fixado no \textit{commit} de prefixo \texttt{40b8c63f75dc}. O pacote foi adquirido antes da coleta, teve seu inventário conferido e não consultou o registro remoto durante as análises. O comando efetivo foi:

\begin{verbatim}
semgrep scan --config /opt/regras-semgrep/python --json \
  --metrics=off --jobs 1 \
  --output /saida/bruto.json.part /entrada
\end{verbatim}

O adaptador preservou identificador da regra, mensagem, caminho, intervalo, CWE, severidade e metadados adicionais. Erros parciais e arquivos ignorados foram mantidos em \texttt{avisos.json}, permitindo distinguir execução concluída com avisos de uma falha do processo.

\section{Execução das Condições de IA (C3 e C4)}\label{sec:execucao-ia-desenvolvimento}

\subsection{Adaptadores de sessão}\label{subsec:adaptadores-sessao-ia}

Cursor e Codex foram integrados por adaptadores diferentes, mas receberam o mesmo código, a mesma instrução-base e o mesmo contrato de saída. Cada chamada correspondeu a uma sessão nova e efêmera. O código Python foi serializado em \texttt{payload-codigo.json}, delimitado como dado não confiável e enviado no \textit{prompt} pela entrada padrão. A pasta de trabalho continha somente o esquema e os arquivos mínimos de configuração do cliente; não continha \texttt{oracle/}, identidade original do projeto, respostas anteriores ou resultados de outras condições.

O Cursor foi executado nativamente no Windows pela versão \texttt{2026.09.26-dd393fe}, com o identificador \texttt{gpt-5.6-luna-medium}. O evento de sistema exibiu \texttt{GPT-5.6 Luna 272K Medium}. Como o cliente informava que o isolamento nativo estava disponível apenas em macOS e Linux, a coleta usou \texttt{--sandbox disabled}, \texttt{--mode=ask} e uma política local que permitia \texttt{Read(**)} e negava \texttt{Write(**)} e \texttt{Shell(*)}. O processo foi iniciado na pasta exclusiva da tentativa, e o código seguiu no \textit{prompt}, não como uma árvore navegável no \textit{workspace}. A auditoria dos eventos de C3 e C5, incluindo os pilotos, encontrou zero chamadas de ferramenta; essa constatação reduz, mas não elimina, a limitação do sandbox nativo indisponível no host.

O Codex foi executado pela \texttt{codex-cli 0.154.0}, solicitando \texttt{gpt-5.6-luna}. O comando aplicou sandbox somente leitura e sessão efêmera, ignorou configurações e regras externas e fixou a política de aprovação como \texttt{never}. Os eventos JSONL dessa versão não expuseram o modelo efetivamente servido. Por isso, \texttt{modelo\_exibido} permaneceu nulo e \texttt{modelo\_verificacao} recebeu \path{nao_exposto_jsonl_codex}; o valor solicitado não foi copiado para o campo observado.

O bloco executado foi o principal, com GPT-5.6 Luna nas quatro condições de IA. O bloco complementar previsto como opcional na metodologia, com Cursor Auto e GPT-5.6 Terra no Codex, não integra o conjunto coletado e não foi combinado com a linha de base.

\subsection{Instrução e contrato de saída}\label{subsec:padronizacao-prompts}

O arquivo \texttt{prompt-relatorio-ia-v1.txt} foi compartilhado por C3--C6 e congelado com SHA-256 de prefixo \texttt{1d3e269aa36d}. A instrução solicitou a inspeção do código fornecido, exigiu uma cadeia de exploração plausível e orientou o agente a considerar validações e controles antes de reportar. Também proibiu executar código, testes, \textit{shell}, rede, busca, MCP, Bandit, Semgrep ou qualquer outro analisador, além de proibir consulta ou reconstrução do \textit{ground truth}.

O contrato de saída, congelado com SHA-256 de prefixo \texttt{6e0b2f29d5f3}, exigiu exclusivamente o \textit{array} JSON descrito na Seção~\ref{subsec:esquema-comum-achados}. O controlador validou localmente a sintaxe, as propriedades, os caminhos, as linhas, os hashes, os metadados de modelo e a ausência de eventos proibidos. O argumento \texttt{--output-schema} não foi aplicado ao Codex, pois essa interface exigia um objeto na raiz e criaria tratamento diferente do Cursor. Em ambos os produtos, a saída bruta foi fechada antes da validação comum.

\section{Construção das Condições Híbridas (C5 e C6)}\label{sec:condicao-hibrida-desenvolvimento}

Após C1 e C2, o script \texttt{gerar-alertas-sast.ps1} selecionou a tentativa concluída autoritativa de cada ferramenta e uniu seus achados sem consultar o \textit{ground truth}. A deduplicação considerou caminho, localização, CWE e, quando necessário, identificação da regra e mensagem. Foram observados 1.416 registros antes da deduplicação e 1.410 alertas por repetição depois da remoção de seis duplicatas.

Os 26 arquivos \texttt{ALVO-NNNN-alertas-sast.json} foram congelados em um \textit{lock} que registra SHA-256, quantidade antes e depois da deduplicação e tentativas C1/C2 de origem. Planejamento e execução de C5 e C6 recusam um arquivo que divirja desse \textit{lock}. Dessa forma, Cursor e Codex receberam exatamente os mesmos bytes de alerta para um mesmo alvo.

O \textit{prompt} comum descreveu os alertas como pistas não adjudicadas. O agente deveria confirmá-los com base no código, ignorar falsos positivos e poderia relatar vulnerabilidades sem alerta correspondente. C5 e C6 foram executadas como sessões novas; não reutilizaram contexto nem respostas de C3 e C4. A única diferença experimental planejada entre a condição isolada e sua correspondente híbrida foi, portanto, a presença dos alertas congelados.

\section{Preservação, Auditoria e Preparação das Métricas}\label{sec:adjudicacao-desenvolvimento}

Em C1 e C2, cada tentativa preservou manifestos de estado, os fluxos de saída e erro em formato binário, saída nativa, versões, dados do processo, avisos e saída normalizada. Em C3--C6, foram preservados \textit{prompt} efetivo, \textit{payload}, alertas quando aplicável, fluxo de eventos, \textit{stderr}, resposta bruta, resposta validada e manifesto. O manifesto terminal referencia todos esses arquivos por SHA-256.

A Tabela~\ref{tab:fechamento-operacional-coleta} apresenta o fechamento operacional. ``Tentativas preservadas'' inclui falhas, retentativas e, em C4, tentativas legadas substituídas na seleção. A tabela não representa acurácia e não emprega o \textit{ground truth}.

\begin{table}[htb]
\ABNTEXfontereduzida
\caption{Fechamento operacional das seis condições}
\label{tab:fechamento-operacional-coleta}
\centering
\begin{tabularx}{\textwidth}{|l|X|r|r|r|}
\hline
\textbf{Condição} & \textbf{Ferramenta e repetições} & \textbf{Previstas} & \textbf{Válidas} & \textbf{Tentativas preservadas} \\
\hline
C1--C2 & Bandit e Semgrep, uma execução por alvo & 52 & 52 & 63 \\
\hline
C3 & Cursor, três repetições por alvo & 78 & 78 & 97 \\
\hline
C4 & Codex, três repetições por alvo & 78 & 78 & 90 \\
\hline
C5 & Cursor com alertas, três repetições por alvo & 78 & 78 & 91 \\
\hline
C6 & Codex com alertas, três repetições por alvo & 78 & 78 & 92 \\
\hline
\textbf{Total} & & \textbf{364} & \textbf{364} & \textbf{433} \\
\hline
\end{tabularx}
\legend{Fonte: Elaborado pelo autor a partir das auditorias de fechamento (2026).}
\end{table}

As auditorias foram calculadas a partir dos registros primários e mantiveram separadas as tentativas válidas e inválidas. Duração e tokens foram registrados por tentativa. No Cursor, \texttt{inputTokens}, \texttt{cacheReadTokens} e \texttt{cacheWriteTokens} compuseram os tokens de entrada, enquanto os componentes originais permaneceram nos eventos. No Codex, os contadores expostos pelo JSONL foram preservados sem estimativa. Essas diferenças de observabilidade devem ser consideradas na análise descritiva do Capítulo~\ref{cap:avaliacao}.

O \textit{oracle} somente pode ser introduzido depois do fechamento da coleta. Os achados normalizados e as respostas validadas constituem as entradas da etapa de correspondência com caminho, CWE e tolerância de linha do RealVuln. Não foram usadas correções manuais para completar localização, CWE ou severidade durante a coleta. TP, FP, FN, TN e as métricas derivadas não aparecem nas auditorias operacionais, evitando que uma síntese sem adjudicação seja confundida com resultado experimental.

\section{Verificação do Artefato}\label{sec:verificacao-artefato-desenvolvimento}

Os adaptadores foram exercitados com \textit{fixtures} artificiais externas ao corpus: saída vazia, achado completo, duplicata, JSON inválido, número não finito, erro parcial, CWE ausente e diferentes formatos de caminho e linha. Os testes do controlador também cobriram políticas de aquisição, links e junções, hashes canônicos, preparação, fila, atomicidade, retomada, timeout, códigos de saída e invariantes dos manifestos.

O \textit{gate} de construção reconstruiu a imagem e executou testes unitários, testes do \textit{wrapper}, integrações Docker, verificação das dependências e sondas de isolamento. A sonda observou usuário sem privilégio, ausência de \textit{capabilities}, \texttt{NoNewPrivs=1}, apenas a interface de rede de \textit{loopback}, raiz somente leitura e ausência de corpus e \textit{oracle} na imagem. Uma fumaça descartável sobre um alvo do RealVuln comprovou o fluxo C1/C2 de ponta a ponta e foi arquivada com \texttt{finalidade=fumaca}, fora da coleta principal.

As interfaces de IA possuem testes próprios que simulam eventos e respostas sem realizar chamadas externas. Antes de cada coleta completa, os quatro cenários também passaram por piloto em \texttt{ALVO-0001}. O planejamento em modo seco validou todas as entradas, configurações e hashes sem consumir API. No fechamento, a repetição desse planejamento informou zero chamadas pendentes em C3--C6.

\section{Configuração Efetivamente Observada}\label{sec:registro-configuracao-desenvolvimento}

A Tabela~\ref{tab:registro-configuracao-efetiva} consolida a configuração observada. Os identificadores de \textit{commit} e SHA-256 aparecem abreviados para leitura; os valores completos estão nos \textit{locks}, manifestos e inventários do artefato.

\begin{table}[htb]
\ABNTEXfontereduzida
\caption{Configuração efetiva da implementação e da coleta}
\label{tab:registro-configuracao-efetiva}
\centering
\begin{tabularx}{\textwidth}{|X|X|X|}
\hline
\textbf{Componente} & \textbf{Valor observado} & \textbf{Fonte de evidência} \\
\hline
RealVuln & v1.0; \textit{commit} \texttt{d98e9fc91273}; 26 alvos & \texttt{fontes.lock.json} e resumo do corpus \\
\hline
Corpus rotulado & 817 entradas; 697 vulnerabilidades; 120 armadilhas & Saída do validador oficial e auditoria do \textit{ground truth} \\
\hline
Ambiente SAST & Python 3.12.13; Bandit 1.9.4; Semgrep CE 1.172.0 & \texttt{requirements.lock}, imagem e \texttt{versoes.json} \\
\hline
Regras Semgrep & \textit{commit} \texttt{40b8c63f75dc}; pacote Python verificado & \texttt{fontes.lock.json} e inventário das regras \\
\hline
Cursor & cliente \texttt{2026.09.26-dd393fe}; \texttt{gpt-5.6-luna-medium}; modelo exibido \texttt{GPT-5.6 Luna 272K Medium} & Eventos e manifestos de C3/C5 \\
\hline
Codex & \texttt{codex-cli 0.154.0}; \texttt{gpt-5.6-luna}; modelo servido não exposto no JSONL & Eventos e manifestos de C4/C6 \\
\hline
Instrução de IA & \textit{prompt} \texttt{1d3e269aa36d}; esquema \texttt{6e0b2f29d5f3} & Configurações dos agentes e manifestos \\
\hline
Alertas híbridos & 1.410 alertas por repetição; seis duplicatas removidas & \textit{Lock} de alertas e auditorias C5/C6 \\
\hline
\end{tabularx}
\legend{Fonte: Elaborado pelo autor (2026).}
\end{table}

O host observado na coleta SAST utilizava Git for Windows 2.46.2, WSL 2.0.14, Docker Desktop 4.38.0 e Docker Engine 27.5.1. Como essas versões estavam abaixo das recomendações registradas e incluíam risco conhecido no Docker Desktop, o artefato preserva uma aceitação explícita de risco residual, limitada ao marco SAST e condicionada aos controles de isolamento. A exceção não foi tratada como autorização genérica para outros experimentos.

\section{Desvios e Limitações Observados}\label{sec:desvios-protocolo}

A execução revelou desvios e limitações que não seriam visíveis em uma descrição puramente prospectiva:

\begin{itemize}
\item Nove tarefas C1 falharam inicialmente por incompatibilidade de formato no adaptador do Bandit. A correção foi aplicada em uma imagem identificada separadamente, as tarefas foram reabertas de forma explícita e concluídas na tentativa 3. Saídas antigas, transições da fila e ambos os identificadores de imagem foram preservados.

\item C4 foi a primeira condição de IA concluída. Setenta unidades selecionadas provinham de uma versão anterior do executor, sem o campo \texttt{modelo\_verificacao} e sem a validação atual de caminho e linha. Oito unidades foram reprocessadas; seis manifestos legados antes marcados como concluídos foram substituídos na seleção, mas não apagados nem reescritos. A auditoria registra 90 tentativas e os dois hashes de configuração históricos.

\item A validação estrita rejeitou respostas de IA que mencionavam arquivo inexistente ou linha fora do arquivo. Foram preservadas 19 falhas de formato em C3, seis em C4, 13 em C5 e 14 em C6. Novas sessões produziram as tentativas seguintes até completar as 78 unidades válidas de cada condição. Uma falha de formato nunca foi convertida em lista vazia ou ausência de achados.

\item No Cursor nativo para Windows, o sandbox de sistema permaneceu desativado. A mitigação combinou área exclusiva, negação de escrita e \textit{shell}, entrada de código por \textit{stdin}, rejeição de eventos proibidos e auditoria posterior. A permissão \texttt{Read(**)} representa capacidade relativa ao \textit{workspace}, não prova de leitura; embora nenhum \texttt{tool\_call} tenha sido observado, essa diferença em relação ao sandbox somente leitura do Codex permanece uma ameaça à validade.

\item A CLI do Codex não expôs no JSONL o modelo efetivamente servido. A implementação preservou essa ausência em vez de inferir que o modelo solicitado foi necessariamente o observado. Em parte do histórico de C4, nem mesmo o campo explicativo atual estava disponível; essa limitação foi mantida nos manifestos legados.

\item Dez alvos e o pacote local de regras Semgrep não puderam ser redistribuídos integralmente. Essa restrição não alterou as entradas da coleta local; afetou apenas a forma de publicação. URLs, revisões, inventários e scripts de reconstrução foram mantidos para permitir auditoria.
\end{itemize}

Ao final, todas as unidades previstas do bloco principal possuíam uma tentativa selecionada e válida, sem tarefas pendentes. A preservação dos desvios, em vez da sua remoção retrospectiva, mantém a rastreabilidade entre configuração, chamada externa, resposta bruta, validação e conjunto escolhido para a avaliação.

%\chapter{Avaliação dos Resultados}\label{cap:avaliacao}
\addcontentsline{toc}{chapter}{Avaliação dos Resultados}

Este capítulo apresenta e discute os resultados obtidos com a execução do protocolo descrito no Capítulo~\ref{cap:metodologia} e relatado no Capítulo~\ref{cap:desenvolvimento}, organizados de acordo com as quatro questões de pesquisa formuladas na Seção~\ref{sec:questoes-pesquisa-metodologia}.

A Seção~\ref{sec:visao-geral-coleta} caracteriza a coleta efetivamente realizada. A Seção~\ref{sec:desempenho-quantitativo} responde à QP1, apresentando o desempenho de detecção por condição, categoria CWE e severidade. A Seção~\ref{sec:qualidade-relatorios} responde à QP2, com os resultados da rubrica qualitativa. A Seção~\ref{sec:efeito-hibrido} responde à QP3, comparando as condições híbridas às condições de IA isolada e SAST isolado. A Seção~\ref{sec:estabilidade-execucoes} e a Seção~\ref{sec:priorizacao-resultados} respondem à QP4, tratando, respectivamente, de estabilidade entre execuções e de priorização por severidade. A Seção~\ref{sec:tempo-tokens-resultados} apresenta a caracterização operacional complementar do tempo de execução e do consumo de tokens, sem vinculá-la a uma questão de pesquisa autônoma. Por fim, a Seção~\ref{sec:discussao-integrada} integra os achados e a Seção~\ref{sec:ameacas-validade-observadas} revisita as ameaças à validade descritas na Seção~\ref{sec:validade-limitacoes} à luz do que foi efetivamente observado.

\begin{quote}
\textbf{Nota de redação:} assim como no Capítulo~\ref{cap:desenvolvimento}, os trechos entre colchetes marcam pontos que dependem dos dados coletados no experimento e serão substituídos por resultados, tabelas e figuras reais. Nenhum valor numérico foi antecipado nesta versão.
\end{quote}

\section{Visão Geral da Coleta e do Conjunto de Dados Analisado}\label{sec:visao-geral-coleta}

[A PREENCHER APÓS A COLETA: número de repositórios, execuções e entradas efetivamente analisadas em cada condição; configuração de modelo usada em cada bloco; confirmação da execução ou não do bloco complementar; e eventuais exclusões ou falhas de execução tratadas pelo modo estrito, conforme a Seção~\ref{sec:universo-amostra}.]

\section{Desempenho Quantitativo de Detecção (QP1)}\label{sec:desempenho-quantitativo}

\subsection{Resultados por condição experimental}\label{subsec:resultados-por-condicao}

[A PREENCHER APÓS A COLETA: precisão, \textit{recall}, F1, F2, F3 e taxa de falsos positivos para C1--C6 no bloco principal com GPT-5.6 Luna, com intervalos de confiança de 95\% obtidos por reamostragem no nível do repositório, conforme a Seção~\ref{subsec:metricas-quantitativas-metodologia}. Se o bloco complementar for executado, apresentar Cursor Auto e Codex com GPT-5.6 Terra em tabela separada, sem combinar essas observações com a linha de base Luna.]

\begin{table}[htb]
\ABNTEXfontereduzida
\caption{\label{tab:resultados-deteccao-condicao}Desempenho de detecção por condição experimental (a preencher)}
\centering
\begin{tabularx}{\textwidth}{|l|X|X|X|X|X|X|}
\hline
\textbf{Condição} & \textbf{Precisão} & \textbf{Recall} & \textbf{F1} & \textbf{F2} & \textbf{F3} & \textbf{Taxa de FP} \\
\hline
C1 & & & & & & \\
\hline
C2 & & & & & & \\
\hline
C3 & & & & & & \\
\hline
C4 & & & & & & \\
\hline
C5 & & & & & & \\
\hline
C6 & & & & & & \\
\hline
\end{tabularx}
\legend{Fonte: Elaborado pelo autor (2026). Valores a preencher após a coleta.}
\end{table}

\subsection{Resultados por categoria CWE}\label{subsec:resultados-por-cwe}

[A PREENCHER APÓS A COLETA: desempenho de cada condição por categoria CWE presente no RealVuln, incluindo casos em que a quantidade de entradas permite apenas interpretação descritiva, conforme previsto na Seção~\ref{subsec:metricas-quantitativas-metodologia}.]

\subsection{Resultados por severidade}\label{subsec:resultados-por-severidade}

[A PREENCHER APÓS A COLETA: desempenho de cada condição por categoria de severidade (baixa, média, alta, crítica) do \textit{ground truth} do RealVuln.]

\section{Qualidade dos Relatórios Verdadeiros-Positivos (QP2)}\label{sec:qualidade-relatorios}

\subsection{Correção técnica, localização e acionabilidade}\label{subsec:resultados-rubrica}

[A PREENCHER APÓS A COLETA: escores médios e distribuição por condição nas três dimensões ordinais da rubrica definida na Tabela~\ref{tab:dimensoes-rubrica}, para o recorte comum e para o recorte estratificado descritos na Seção~\ref{subsec:metricas-qualitativas-metodologia}. Se houver bloco complementar, apresentar sua amostra e seus resultados separadamente do bloco Luna.]

\subsection{Clareza e legibilidade}\label{subsec:resultados-clareza}

[A PREENCHER APÓS A COLETA: resultados do checklist estrutural de clareza (Tabela~\ref{tab:checklist-clareza}) e do Índice de Legibilidade de Flesch adaptado ao português, por condição.]

\subsection{Concordância entre avaliadores}\label{subsec:resultados-concordancia}

[A PREENCHER APÓS A COLETA: Kappa ponderado de Cohen para as dimensões ordinais e percentual de concordância e Kappa não ponderado para os itens binários de clareza, conforme a Seção~\ref{subsec:metricas-qualitativas-metodologia}, com discussão de eventuais divergências registradas entre avaliadores.]

\section{Efeito da Combinação SAST + IA (QP3)}\label{sec:efeito-hibrido}

[A PREENCHER APÓS A COLETA: comparação pareada entre C3/C4 (IA isolada) e C5/C6 (híbrido correspondente), com intervalos de confiança pareados, discutindo se o fornecimento dos alertas de Bandit e Semgrep aumentou ou reduziu a detecção e a qualidade dos relatórios, à luz do alerta metodológico de \cite{Croft2021RuleLearningSAST} citado na Seção~\ref{subsubsec:sast-vantagens-limitacoes}. Realizar as comparações dentro de cada bloco de modelos, sem parear Luna com Auto ou Terra.]

\section{Estabilidade entre Execuções (QP4)}\label{sec:estabilidade-execucoes}

[A PREENCHER APÓS A COLETA: média, amplitude e desvio das métricas entre as três repetições de C3--C6, e frequência de detecção (zero, uma, duas ou três execuções) por entrada vulnerável, conforme a Seção~\ref{subsec:testes-ensaios}. Calcular a estabilidade separadamente para o bloco Luna e para o bloco complementar, caso realizado.]

\section{Priorização e Indicação de Severidade (QP4)}\label{sec:priorizacao-resultados}

[A PREENCHER APÓS A COLETA: cobertura de priorização, concordância exata, concordância com tolerância de uma categoria e Kappa ponderado entre a severidade indicada por cada condição e a categoria de referência do RealVuln, conforme a Seção~\ref{subsec:priorizacao-metodologia}.]

\section{Tempo de Execução e Consumo de Tokens}\label{sec:tempo-tokens-resultados}

[A PREENCHER APÓS A COLETA: mediana, intervalo interquartil e valores mínimos e máximos do tempo por condição e por bloco de modelos. Para C3--C6, apresentar também tokens de entrada, de saída e totais por execução e por repositório, além de categorias adicionais informadas pelos produtos. Indicar explicitamente a proporção de execuções com contadores disponíveis e evitar comparações diretas quando as definições ou os tokenizadores não forem equivalentes, conforme a Seção~\ref{subsec:tempo-tokens-metodologia}.]

\section{Discussão Integrada}\label{sec:discussao-integrada}

[A PREENCHER APÓS A COLETA: leitura conjunta dos resultados de detecção, qualidade, estabilidade e priorização, acompanhada da caracterização de tempo e consumo de tokens, evitando reduzir a comparação a uma única métrica, conforme a advertência de validade de construto da Seção~\ref{sec:validade-limitacoes}. Discussão explícita de como os achados dialogam com os estudos do Capítulo~\ref{cap:estadoarte}.]

\section{Ameaças à Validade Observadas na Coleta}\label{sec:ameacas-validade-observadas}

[A PREENCHER APÓS A COLETA: revisão da Tabela~\ref{tab:ameacas-validade}, indicando quais ameaças se manifestaram de fato durante a coleta, como foram tratadas e quais limitações adicionais só puderam ser identificadas após a execução do experimento.]


% PARTE
%\part{Desenvolvimento}
%\chapter{Sistema Proposto}\label{cap:proposta}

Esse trabalho propõe um sistema de... 


\section{Primeira Parte do Sistema Proposto}

\lipsum[67]

\section{Considerações Finais}

\lipsum[68]


% PARTE
%\part{Parte Final}
%\chapter{Avaliação dos Resultados}\label{cap:avaliacao}
\addcontentsline{toc}{chapter}{Avaliação dos Resultados}

Este capítulo apresenta e discute os resultados obtidos com a execução do protocolo descrito no Capítulo~\ref{cap:metodologia} e relatado no Capítulo~\ref{cap:desenvolvimento}, organizados de acordo com as quatro questões de pesquisa formuladas na Seção~\ref{sec:questoes-pesquisa-metodologia}.

A Seção~\ref{sec:visao-geral-coleta} caracteriza a coleta efetivamente realizada. A Seção~\ref{sec:desempenho-quantitativo} responde à QP1, apresentando o desempenho de detecção por condição, categoria CWE e severidade. A Seção~\ref{sec:qualidade-relatorios} responde à QP2, com os resultados da rubrica qualitativa. A Seção~\ref{sec:efeito-hibrido} responde à QP3, comparando as condições híbridas às condições de IA isolada e SAST isolado. A Seção~\ref{sec:estabilidade-execucoes} e a Seção~\ref{sec:priorizacao-resultados} respondem à QP4, tratando, respectivamente, de estabilidade entre execuções e de priorização por severidade. A Seção~\ref{sec:tempo-tokens-resultados} apresenta a caracterização operacional complementar do tempo de execução e do consumo de tokens, sem vinculá-la a uma questão de pesquisa autônoma. Por fim, a Seção~\ref{sec:discussao-integrada} integra os achados e a Seção~\ref{sec:ameacas-validade-observadas} revisita as ameaças à validade descritas na Seção~\ref{sec:validade-limitacoes} à luz do que foi efetivamente observado.

\begin{quote}
\textbf{Nota de redação:} assim como no Capítulo~\ref{cap:desenvolvimento}, os trechos entre colchetes marcam pontos que dependem dos dados coletados no experimento e serão substituídos por resultados, tabelas e figuras reais. Nenhum valor numérico foi antecipado nesta versão.
\end{quote}

\section{Visão Geral da Coleta e do Conjunto de Dados Analisado}\label{sec:visao-geral-coleta}

[A PREENCHER APÓS A COLETA: número de repositórios, execuções e entradas efetivamente analisadas em cada condição; configuração de modelo usada em cada bloco; confirmação da execução ou não do bloco complementar; e eventuais exclusões ou falhas de execução tratadas pelo modo estrito, conforme a Seção~\ref{sec:universo-amostra}.]

\section{Desempenho Quantitativo de Detecção (QP1)}\label{sec:desempenho-quantitativo}

\subsection{Resultados por condição experimental}\label{subsec:resultados-por-condicao}

[A PREENCHER APÓS A COLETA: precisão, \textit{recall}, F1, F2, F3 e taxa de falsos positivos para C1--C6 no bloco principal com GPT-5.6 Luna, com intervalos de confiança de 95\% obtidos por reamostragem no nível do repositório, conforme a Seção~\ref{subsec:metricas-quantitativas-metodologia}. Se o bloco complementar for executado, apresentar Cursor Auto e Codex com GPT-5.6 Terra em tabela separada, sem combinar essas observações com a linha de base Luna.]

\begin{table}[htb]
\ABNTEXfontereduzida
\caption{\label{tab:resultados-deteccao-condicao}Desempenho de detecção por condição experimental (a preencher)}
\centering
\begin{tabularx}{\textwidth}{|l|X|X|X|X|X|X|}
\hline
\textbf{Condição} & \textbf{Precisão} & \textbf{Recall} & \textbf{F1} & \textbf{F2} & \textbf{F3} & \textbf{Taxa de FP} \\
\hline
C1 & & & & & & \\
\hline
C2 & & & & & & \\
\hline
C3 & & & & & & \\
\hline
C4 & & & & & & \\
\hline
C5 & & & & & & \\
\hline
C6 & & & & & & \\
\hline
\end{tabularx}
\legend{Fonte: Elaborado pelo autor (2026). Valores a preencher após a coleta.}
\end{table}

\subsection{Resultados por categoria CWE}\label{subsec:resultados-por-cwe}

[A PREENCHER APÓS A COLETA: desempenho de cada condição por categoria CWE presente no RealVuln, incluindo casos em que a quantidade de entradas permite apenas interpretação descritiva, conforme previsto na Seção~\ref{subsec:metricas-quantitativas-metodologia}.]

\subsection{Resultados por severidade}\label{subsec:resultados-por-severidade}

[A PREENCHER APÓS A COLETA: desempenho de cada condição por categoria de severidade (baixa, média, alta, crítica) do \textit{ground truth} do RealVuln.]

\section{Qualidade dos Relatórios Verdadeiros-Positivos (QP2)}\label{sec:qualidade-relatorios}

\subsection{Correção técnica, localização e acionabilidade}\label{subsec:resultados-rubrica}

[A PREENCHER APÓS A COLETA: escores médios e distribuição por condição nas três dimensões ordinais da rubrica definida na Tabela~\ref{tab:dimensoes-rubrica}, para o recorte comum e para o recorte estratificado descritos na Seção~\ref{subsec:metricas-qualitativas-metodologia}. Se houver bloco complementar, apresentar sua amostra e seus resultados separadamente do bloco Luna.]

\subsection{Clareza e legibilidade}\label{subsec:resultados-clareza}

[A PREENCHER APÓS A COLETA: resultados do checklist estrutural de clareza (Tabela~\ref{tab:checklist-clareza}) e do Índice de Legibilidade de Flesch adaptado ao português, por condição.]

\subsection{Concordância entre avaliadores}\label{subsec:resultados-concordancia}

[A PREENCHER APÓS A COLETA: Kappa ponderado de Cohen para as dimensões ordinais e percentual de concordância e Kappa não ponderado para os itens binários de clareza, conforme a Seção~\ref{subsec:metricas-qualitativas-metodologia}, com discussão de eventuais divergências registradas entre avaliadores.]

\section{Efeito da Combinação SAST + IA (QP3)}\label{sec:efeito-hibrido}

[A PREENCHER APÓS A COLETA: comparação pareada entre C3/C4 (IA isolada) e C5/C6 (híbrido correspondente), com intervalos de confiança pareados, discutindo se o fornecimento dos alertas de Bandit e Semgrep aumentou ou reduziu a detecção e a qualidade dos relatórios, à luz do alerta metodológico de \cite{Croft2021RuleLearningSAST} citado na Seção~\ref{subsubsec:sast-vantagens-limitacoes}. Realizar as comparações dentro de cada bloco de modelos, sem parear Luna com Auto ou Terra.]

\section{Estabilidade entre Execuções (QP4)}\label{sec:estabilidade-execucoes}

[A PREENCHER APÓS A COLETA: média, amplitude e desvio das métricas entre as três repetições de C3--C6, e frequência de detecção (zero, uma, duas ou três execuções) por entrada vulnerável, conforme a Seção~\ref{subsec:testes-ensaios}. Calcular a estabilidade separadamente para o bloco Luna e para o bloco complementar, caso realizado.]

\section{Priorização e Indicação de Severidade (QP4)}\label{sec:priorizacao-resultados}

[A PREENCHER APÓS A COLETA: cobertura de priorização, concordância exata, concordância com tolerância de uma categoria e Kappa ponderado entre a severidade indicada por cada condição e a categoria de referência do RealVuln, conforme a Seção~\ref{subsec:priorizacao-metodologia}.]

\section{Tempo de Execução e Consumo de Tokens}\label{sec:tempo-tokens-resultados}

[A PREENCHER APÓS A COLETA: mediana, intervalo interquartil e valores mínimos e máximos do tempo por condição e por bloco de modelos. Para C3--C6, apresentar também tokens de entrada, de saída e totais por execução e por repositório, além de categorias adicionais informadas pelos produtos. Indicar explicitamente a proporção de execuções com contadores disponíveis e evitar comparações diretas quando as definições ou os tokenizadores não forem equivalentes, conforme a Seção~\ref{subsec:tempo-tokens-metodologia}.]

\section{Discussão Integrada}\label{sec:discussao-integrada}

[A PREENCHER APÓS A COLETA: leitura conjunta dos resultados de detecção, qualidade, estabilidade e priorização, acompanhada da caracterização de tempo e consumo de tokens, evitando reduzir a comparação a uma única métrica, conforme a advertência de validade de construto da Seção~\ref{sec:validade-limitacoes}. Discussão explícita de como os achados dialogam com os estudos do Capítulo~\ref{cap:estadoarte}.]

\section{Ameaças à Validade Observadas na Coleta}\label{sec:ameacas-validade-observadas}

[A PREENCHER APÓS A COLETA: revisão da Tabela~\ref{tab:ameacas-validade}, indicando quais ameaças se manifestaram de fato durante a coleta, como foram tratadas e quais limitações adicionais só puderam ser identificadas após a execução do experimento.]

%\chapter{Conclusão}\label{cap:conclusao}
\addcontentsline{toc}{chapter}{Conclusão}

Este capítulo encerra o trabalho retomando, à luz dos resultados discutidos no Capítulo~\ref{cap:avaliacao}, o problema de pesquisa e as questões formuladas na Seção~\ref{sec:questoes-pesquisa-metodologia}. A Seção~\ref{sec:sintese-trabalho} sintetiza o percurso do trabalho. A Seção~\ref{sec:resposta-questoes-pesquisa} responde diretamente a cada questão de pesquisa. A Seção~\ref{sec:contribuicoes} indica as contribuições do estudo. A Seção~\ref{sec:limitacoes-conclusao} consolida as limitações reconhecidas. A Seção~\ref{sec:trabalhos-futuros} aponta direções de continuidade.

\begin{quote}
\textbf{Nota de redação:} como nos capítulos anteriores, os trechos entre colchetes marcam pontos que só podem ser redigidos após a coleta e a análise dos resultados descritos no Capítulo~\ref{cap:avaliacao}. Nenhuma conclusão factual foi antecipada nesta versão.
\end{quote}

\section{Síntese do Trabalho}\label{sec:sintese-trabalho}

[A PREENCHER APÓS A COLETA: retomada breve do problema de pesquisa --- comparação entre SAST tradicional, agentes empresariais de IA e abordagem híbrida SAST + IA na detecção de vulnerabilidades em Python --- e do percurso metodológico e experimental adotado, sem repetir integralmente o conteúdo dos capítulos anteriores.]

\section{Resposta às Questões de Pesquisa}\label{sec:resposta-questoes-pesquisa}

[A PREENCHER APÓS A COLETA: resposta objetiva a cada uma das questões QP1--QP5 definidas na Seção~\ref{sec:questoes-pesquisa-metodologia}, com remissão direta às seções correspondentes do Capítulo~\ref{cap:avaliacao} onde cada resultado foi apresentado.]

\begin{itemize}
\item \textbf{QP1} --- desempenho comparativo de detecção: [A PREENCHER, ver Seção~\ref{sec:desempenho-quantitativo}].
\item \textbf{QP2} --- qualidade dos relatórios verdadeiros-positivos: [A PREENCHER, ver Seção~\ref{sec:qualidade-relatorios}].
\item \textbf{QP3} --- efeito de fornecer alertas SAST aos agentes de IA: [A PREENCHER, ver Seção~\ref{sec:efeito-hibrido}].
\item \textbf{QP4} --- estabilidade entre execuções e adequação da priorização: [A PREENCHER, ver Seções~\ref{sec:estabilidade-execucoes} e~\ref{sec:priorizacao-resultados}].
\item \textbf{QP5} --- tempo de execução e custo observável: [A PREENCHER, ver Seção~\ref{sec:tempo-custo-resultados}].
\end{itemize}

\section{Contribuições}\label{sec:contribuicoes}

[A PREENCHER APÓS A COLETA: contribuições efetivamente entregues pelo trabalho, por exemplo a comparação tripla SAST--IA--híbrido em Python identificada como lacuna na Seção~\ref{sec:lacunas-literatura}, o \textit{harness} e os artefatos de reprodução descritos no Capítulo~\ref{cap:desenvolvimento}, e a rubrica qualitativa aplicada aos relatórios das seis condições.]

\section{Limitações}\label{sec:limitacoes-conclusao}

[A PREENCHER APÓS A COLETA: consolidação das limitações já antecipadas na Seção~\ref{sec:validade-limitacoes} e das limitações adicionais identificadas apenas durante a execução, registradas na Seção~\ref{sec:ameacas-validade-observadas}, incluindo a delimitação da generalização das conclusões ao RealVuln v1.0 e às versões, regras, \textit{prompts} e ferramentas efetivamente avaliadas.]

\section{Trabalhos Futuros}\label{sec:trabalhos-futuros}

[A PREENCHER APÓS A COLETA: direções de continuidade sugeridas pelos resultados obtidos --- por exemplo, extensão a outros benchmarks ou linguagens, avaliação de outras ferramentas empresariais de IA, refinamento dos desenhos híbridos discutidos na Seção~\ref{sec:abordagens-hibridas}, ou aprofundamento da avaliação de priorização e severidade.]


% ----------------------------------------------------------
% ELEMENTOS PÓS-TEXTUAIS (Referências, Glossário, Apêndices)
% ----------------------------------------------------------
\postextual

% Referências bibliográficas
\nocite{owasp_sast,bandit_docs,semgrep_ce,cursor_home,cursor_agents,openai_codex_cloud,openai_codex_cli,openai_codex_pricing_2026,stackoverflow2025,robbes2026agentic,becker2025productivity,realvuln2026,pygoat_owasp}
\bibliography{bibliografia}

% Glossário (Consulte o manual)
%\glossary

% Apêndices
%% ----------------------------------------------------------
% Apêndices
% ----------------------------------------------------------

% ---
% Inicia os apêndices
% ---
\begin{apendicesenv}

% Imprime uma página indicando o início dos apêndices
\partapendices

% ----------------------------------------------------------
\chapter{Primeiro Apêncice}
% ----------------------------------------------------------

\lipsum[50] % Texto qualquer. REMOVER!!

% ----------------------------------------------------------
\chapter{Segundo apêndice com título tão grande quanto se queira porque ele já faz a quebra de linha da coisa toda}
% ----------------------------------------------------------
\lipsum[51-53] % Texto qualquer. REMOVER!!

\end{apendicesenv}
% ---

% Anexos
%% ----------------------------------------------------------
% Apêndices
% ----------------------------------------------------------

% ---
% Inicia os anexos
% ---
\begin{anexosenv}

% Imprime uma página indicando o início dos anexos
\partanexos

% ---
\chapter{Nome do Primeiro Anexo}
% ---
\lipsum[30] % Texto qualquer. REMOVER!!

% ---
\chapter{Nome de Outro Anexo}
% ---

\lipsum[32] % Texto qualquer. REMOVER!!

\end{anexosenv}

% Índice remissivo (Consultar manual)
\phantompart
\printindex

\end{document}
